% Generated human-review packet template.  Source records, Lean statements,
% and raw-source-to-Spec screening fields are substituted by
% scripts/review_dashboard_packet.py.
% Do not hand-edit generated packets; annotate the PDF or reconcile notes into
% the paper-local human-review log.
\documentclass[10pt]{article}
\usepackage[margin=0.43in,footskip=0.2in]{geometry}
\usepackage{fontspec}
% Use a Unicode-complete main face as well as a Unicode-complete code face.
% Reviewer reasons and source labels can contain exact Greek symbols outside
% verbatim blocks; silently dropping one would corrupt the review packet.
\setmainfont{DejaVu Serif}
% The broad DejaVu Sans Unicode coverage preserves Lean's mathematical glyphs
% (including U+2216 set difference and superscript identifiers) in verbatim
% review blocks.  Its proportional metrics are preferable to silently missing
% exact semantic text from a narrower monospace font.
\setmonofont{DejaVu Sans}
\usepackage{hyperref}
\usepackage{amssymb}
\usepackage{fvextra}
\usepackage{enumitem}
\setlength{\parindent}{0pt}
\setlength{\parskip}{0.15em}
\linespread{0.92}
\hypersetup{colorlinks=true,linkcolor=blue,urlcolor=blue,breaklinks=true}
\DefineVerbatimEnvironment{ReviewVerbatim}{Verbatim}{breaklines=true,breakanywhere=true,fontsize=\scriptsize}
\newcommand{\reviewerbox}{%
  \par\noindent\fbox{\parbox{0.96\linewidth}{\rule{0pt}{0.38in}}}\par
}
\newcommand{\reviewmatch}[1]{%
  \par\noindent\textbf{Reviewer check:} $\square$\; #1\par
  \noindent\emph{If left unchecked, explain the discrepancy or uncertainty below.}\par
}

\begin{document}
\fontsize{9}{10}\selectfont

\begin{center}
{\LARGE Human review packet: Galarza26RelativeCredit}\\[0.35em]
\small Generated from byte-pinned source excerpts, the expanded PaperInterface
specifications, and hash-bound raw-source-to-Spec screening records.
\end{center}

\noindent\textbf{Paper:} Galarza26RelativeCredit\\
\textbf{Paper id:} \texttt{Galarza26RelativeCredit}\\
\textbf{Source version:} \parbox[t]{0.76\linewidth}{\raggedright Versión para entrega, 22 de noviembre de 2026, commit 25f0b0c7df040e2b04d78064b30949492cd4472f}\\
\textbf{Source record:} \texttt{audit/paper\_statement\_map.json}\\
\textbf{Rows:} 3\\
\textbf{Generated:} 2026-11-22\\
\textbf{Source URL:} \href{https://github.com/fgalarzac/ai-project/blob/25f0b0c7df040e2b04d78064b30949492cd4472f/paper/paper.pdf}{official source}





\section*{How to use this packet}
For each source claim, compare the exact source excerpts with the expanded
PaperInterface specification. Mark up this PDF or fill the blank reviewer box in the TeX source;
return the annotated file for reconciliation into the paper-local human-review
log.

\section*{Open the interactive dashboard}
The interactive dashboard is an optional alternative to using this PDF; it
presents the same review surface rather than an additional review requirement.
After forking and cloning (or pulling) AppliedModelingLib, run the following from the
repository root. The cache command is needed only the first time in a new clone.
\begin{ReviewVerbatim}
cd <your-repository-checkout>
lake exe cache get
python3 scripts/review_dashboard.py --paper Galarza26RelativeCredit --serve
\end{ReviewVerbatim}
Open \url{http://127.0.0.1:8765/} in a browser and keep that terminal running
while reviewing. The dashboard records saved annotations in this paper's local
reviewer-owned review record.

\section*{Contents}
\begin{itemize}[leftmargin=1.2em]
\item \hyperlink{governing-declaration-section-5ef5ef0364b6939c}{Governing Lean declarations (9; supporting context)}
\item \hyperlink{paper-prerequisite-section-5ef5ef0364b6939c}{Paper-specific semantic prerequisites (18; not paper claims)}
\begin{itemize}[leftmargin=1.2em]
\item \hyperlink{paper-prerequisite-6956aaccf5fffc8f}{Galarza26\allowbreak{}Relative\allowbreak{}Credit.\allowbreak{}premium}
\item \hyperlink{paper-prerequisite-e8657fe6a37f000a}{Galarza26\allowbreak{}Relative\allowbreak{}Credit.\allowbreak{}prior\allowbreak{}Risk}
\item \hyperlink{paper-prerequisite-3f518444e1a02f7c}{Galarza26\allowbreak{}Relative\allowbreak{}Credit.\allowbreak{}screening\allowbreak{}Cost}
\item \hyperlink{paper-prerequisite-646bdd87bd4c2b85}{Galarza26\allowbreak{}Relative\allowbreak{}Credit.\allowbreak{}bank\allowbreak{}Profit}
\item \hyperlink{paper-prerequisite-b0571b2c4a78eab5}{Galarza26\allowbreak{}Relative\allowbreak{}Credit.\allowbreak{}cohort\allowbreak{}Weight}
\item \hyperlink{paper-prerequisite-3564270d1ab44650}{Galarza26\allowbreak{}Relative\allowbreak{}Credit.\allowbreak{}feasible\allowbreak{}Efforts}
\item \hyperlink{paper-prerequisite-a6fae381e2d9a684}{Galarza26\allowbreak{}Relative\allowbreak{}Credit.\allowbreak{}conditioned\allowbreak{}Continuation\allowbreak{}Problem}
\item \hyperlink{paper-prerequisite-0bf7e0c122a51179}{Galarza26\allowbreak{}Relative\allowbreak{}Credit.\allowbreak{}conditioned\allowbreak{}Payoff}
\item \hyperlink{paper-prerequisite-edcf6b4879cafd1b}{Galarza26\allowbreak{}Relative\allowbreak{}Credit.\allowbreak{}screening\allowbreak{}Strength}
\item \hyperlink{paper-prerequisite-29fcb63a4a0578b3}{Galarza26\allowbreak{}Relative\allowbreak{}Credit.\allowbreak{}screening\allowbreak{}Advantage}
\item \hyperlink{paper-prerequisite-4f59b48de0be00ac}{Galarza26\allowbreak{}Relative\allowbreak{}Credit.\allowbreak{}crossing\allowbreak{}Conditions}
\item \hyperlink{paper-prerequisite-527f5e3e1ca86ece}{Galarza26\allowbreak{}Relative\allowbreak{}Credit.\allowbreak{}impaired\allowbreak{}Balance}
\item \hyperlink{paper-prerequisite-3b6383e0216b3c72}{Galarza26\allowbreak{}Relative\allowbreak{}Credit.\allowbreak{}relative\allowbreak{}Price}
\item \hyperlink{paper-prerequisite-38fc4c8d2be8054c}{Galarza26\allowbreak{}Relative\allowbreak{}Credit.\allowbreak{}relative\allowbreak{}Price\_\allowbreak{}eq\_\allowbreak{}rate\_\allowbreak{}gap}
\item \hyperlink{paper-prerequisite-e9553b2c455b69dd}{Galarza26\allowbreak{}Relative\allowbreak{}Credit.\allowbreak{}screened\allowbreak{}Risk}
\item \hyperlink{paper-prerequisite-ab0e2f9b84884bc6}{Galarza26\allowbreak{}Relative\allowbreak{}Credit.\allowbreak{}selection\allowbreak{}And\allowbreak{}Domain}
\item \hyperlink{paper-prerequisite-de4cba6057d51d83}{Galarza26\allowbreak{}Relative\allowbreak{}Credit.\allowbreak{}stock\allowbreak{}Accounting\allowbreak{}Conditions}
\item \hyperlink{paper-prerequisite-e5f9b7ce782f743a}{Galarza26\allowbreak{}Relative\allowbreak{}Credit.\allowbreak{}technology\allowbreak{}Capacity\allowbreak{}Cost}
\end{itemize}
\item \hyperlink{source-section-f10b5f68e4acf3dc}{Source claims (3)}
\begin{itemize}[leftmargin=1.2em]
\item \hyperlink{source-claim-71d1c57a6546a5a7}{Proposición 1 (Solución global y pendientes por régimen)}
\item \hyperlink{source-claim-f0d8e27dfc3b3ee3}{Proposición 2 (Umbral condicional)}
\item \hyperlink{source-claim-194f750a63a93998}{Proposición 3 (Transmisión y composición)}
\end{itemize}
\end{itemize}

\clearpage
\hypertarget{governing-declaration-section-5ef5ef0364b6939c}{}
\section*{Governing Lean declarations}
These exact graph-bound declaration bodies are retained by the displayed specifications or reviewed prerequisites. They are supporting context and do not add source-result rows or semantic judgments.
\hypertarget{governing-declaration-34d15e1983b4c930}{}
\subsection*{\small\ttfamily\raggedright Galarza26\allowbreak{}Relative\allowbreak{}Credit.\allowbreak{}Continuation\allowbreak{}Problem}
\textbf{Declaration location:} paper-local
\paragraph{Lean declaration body}
\begin{ReviewVerbatim}
field feasible:
Galarza26RelativeCredit.ContinuationProblem → Set ℝ

field payoff:
Galarza26RelativeCredit.ContinuationProblem → ℝ → ℝ
\end{ReviewVerbatim}


\hypertarget{governing-declaration-042b6300ffc11e3e}{}
\subsection*{\small\ttfamily\raggedright Galarza26\allowbreak{}Relative\allowbreak{}Credit.\allowbreak{}Continuation\allowbreak{}Problem.\allowbreak{}Is\allowbreak{}Maximizer}
\textbf{Declaration location:} paper-local
\paragraph{Lean declaration body}
\begin{ReviewVerbatim}
type:
Galarza26RelativeCredit.ContinuationProblem → ℝ → Prop

value:
fun (problem : Galarza26RelativeCredit.ContinuationProblem) (e : ℝ) =>
  e ∈ problem.feasible ∧ ∀ x ∈ problem.feasible, problem.payoff x ≤ problem.payoff e
\end{ReviewVerbatim}

\paragraph{Direct retained dependencies}
\begin{itemize}\raggedright
\item \hyperlink{governing-declaration-34d15e1983b4c930}{Galarza26\allowbreak{}Relative\allowbreak{}Credit.\allowbreak{}Continuation\allowbreak{}Problem}
\end{itemize}
\hypertarget{governing-declaration-ac78c6a99f925061}{}
\subsection*{\small\ttfamily\raggedright Galarza26\allowbreak{}Relative\allowbreak{}Credit.\allowbreak{}exponential\allowbreak{}Balance}
\textbf{Declaration location:} paper-local
\paragraph{Lean declaration body}
\begin{ReviewVerbatim}
type:
ℝ → ℝ → ℝ → ℝ

value:
fun (B1 epsilon D : ℝ) => B1 * Real.exp (-epsilon * D)
\end{ReviewVerbatim}


\hypertarget{governing-declaration-c80cb553f94a2c2b}{}
\subsection*{\small\ttfamily\raggedright Galarza26\allowbreak{}Relative\allowbreak{}Credit.\allowbreak{}portfolio\allowbreak{}Ratio}
\textbf{Declaration location:} paper-local
\paragraph{Lean declaration body}
\begin{ReviewVerbatim}
type:
ℝ → ℝ → (ℝ → ℝ) → (ℝ → ℝ) → ℝ → ℝ

value:
fun (B0 m0 : ℝ) (B q : ℝ → ℝ) (D : ℝ) => Galarza26RelativeCredit.impairedBalance B0 m0 B q D / (B0 + B D)
\end{ReviewVerbatim}

\paragraph{Direct retained dependencies}
\begin{itemize}\raggedright
\item \hyperlink{paper-prerequisite-527f5e3e1ca86ece}{Galarza26\allowbreak{}Relative\allowbreak{}Credit.\allowbreak{}impaired\allowbreak{}Balance}
\end{itemize}
\hypertarget{governing-declaration-6e5d4449ad412574}{}
\subsection*{\small\ttfamily\raggedright Galarza26\allowbreak{}Relative\allowbreak{}Credit.\allowbreak{}capacity\allowbreak{}Regime}
\textbf{Declaration location:} paper-local
\paragraph{Lean declaration body}
\begin{ReviewVerbatim}
type:
ℝ → ℝ → ℝ → ℝ → ℝ → ℝ → ℝ → Prop

value:
fun (qbar a g c Rbar E D : ℝ) =>
  E < g * Galarza26RelativeCredit.premium Rbar D / c ∧ E < Galarza26RelativeCredit.priorRisk qbar a D / g
\end{ReviewVerbatim}

\paragraph{Direct retained dependencies}
\begin{itemize}\raggedright
\item \hyperlink{paper-prerequisite-6956aaccf5fffc8f}{Galarza26\allowbreak{}Relative\allowbreak{}Credit.\allowbreak{}premium}
\item \hyperlink{paper-prerequisite-e8657fe6a37f000a}{Galarza26\allowbreak{}Relative\allowbreak{}Credit.\allowbreak{}prior\allowbreak{}Risk}
\end{itemize}
\hypertarget{governing-declaration-65adc159d5a1934d}{}
\subsection*{\small\ttfamily\raggedright Galarza26\allowbreak{}Relative\allowbreak{}Credit.\allowbreak{}floor\allowbreak{}Regime}
\textbf{Declaration location:} paper-local
\paragraph{Lean declaration body}
\begin{ReviewVerbatim}
type:
ℝ → ℝ → ℝ → ℝ → ℝ → ℝ → ℝ → Prop

value:
fun (qbar a g c Rbar E D : ℝ) =>
  Galarza26RelativeCredit.priorRisk qbar a D / g < g * Galarza26RelativeCredit.premium Rbar D / c ∧
    Galarza26RelativeCredit.priorRisk qbar a D / g < E
\end{ReviewVerbatim}

\paragraph{Direct retained dependencies}
\begin{itemize}\raggedright
\item \hyperlink{paper-prerequisite-6956aaccf5fffc8f}{Galarza26\allowbreak{}Relative\allowbreak{}Credit.\allowbreak{}premium}
\item \hyperlink{paper-prerequisite-e8657fe6a37f000a}{Galarza26\allowbreak{}Relative\allowbreak{}Credit.\allowbreak{}prior\allowbreak{}Risk}
\end{itemize}
\hypertarget{governing-declaration-5f246d7c1bbfcd61}{}
\subsection*{\small\ttfamily\raggedright Galarza26\allowbreak{}Relative\allowbreak{}Credit.\allowbreak{}interior\allowbreak{}Regime}
\textbf{Declaration location:} paper-local
\paragraph{Lean declaration body}
\begin{ReviewVerbatim}
type:
ℝ → ℝ → ℝ → ℝ → ℝ → ℝ → ℝ → Prop

value:
fun (qbar a g c Rbar E D : ℝ) =>
  g * Galarza26RelativeCredit.premium Rbar D / c < E ∧
    g * Galarza26RelativeCredit.premium Rbar D / c < Galarza26RelativeCredit.priorRisk qbar a D / g
\end{ReviewVerbatim}

\paragraph{Direct retained dependencies}
\begin{itemize}\raggedright
\item \hyperlink{paper-prerequisite-6956aaccf5fffc8f}{Galarza26\allowbreak{}Relative\allowbreak{}Credit.\allowbreak{}premium}
\item \hyperlink{paper-prerequisite-e8657fe6a37f000a}{Galarza26\allowbreak{}Relative\allowbreak{}Credit.\allowbreak{}prior\allowbreak{}Risk}
\end{itemize}
\hypertarget{governing-declaration-41e961e1975be029}{}
\subsection*{\small\ttfamily\raggedright Galarza26\allowbreak{}Relative\allowbreak{}Credit.\allowbreak{}optimal\allowbreak{}Effort}
\textbf{Declaration location:} paper-local
\paragraph{Lean declaration body}
\begin{ReviewVerbatim}
type:
ℝ → ℝ → ℝ → ℝ → ℝ → ℝ → ℝ → ℝ

value:
fun (qbar a g c Rbar E D : ℝ) =>
  min (g * Galarza26RelativeCredit.premium Rbar D / c) (min E (Galarza26RelativeCredit.priorRisk qbar a D / g))
\end{ReviewVerbatim}

\paragraph{Direct retained dependencies}
\begin{itemize}\raggedright
\item \hyperlink{paper-prerequisite-6956aaccf5fffc8f}{Galarza26\allowbreak{}Relative\allowbreak{}Credit.\allowbreak{}premium}
\item \hyperlink{paper-prerequisite-e8657fe6a37f000a}{Galarza26\allowbreak{}Relative\allowbreak{}Credit.\allowbreak{}prior\allowbreak{}Risk}
\end{itemize}
\hypertarget{governing-declaration-e139cf967ddac229}{}
\subsection*{\small\ttfamily\raggedright Galarza26\allowbreak{}Relative\allowbreak{}Credit.\allowbreak{}optimal\allowbreak{}Risk}
\textbf{Declaration location:} paper-local
\paragraph{Lean declaration body}
\begin{ReviewVerbatim}
type:
ℝ → ℝ → ℝ → ℝ → ℝ → ℝ → ℝ → ℝ

value:
fun (qbar a g c Rbar E D : ℝ) =>
  max 0 (Galarza26RelativeCredit.priorRisk qbar a D - g * min E (g * Galarza26RelativeCredit.premium Rbar D / c))
\end{ReviewVerbatim}

\paragraph{Direct retained dependencies}
\begin{itemize}\raggedright
\item \hyperlink{paper-prerequisite-6956aaccf5fffc8f}{Galarza26\allowbreak{}Relative\allowbreak{}Credit.\allowbreak{}premium}
\item \hyperlink{paper-prerequisite-e8657fe6a37f000a}{Galarza26\allowbreak{}Relative\allowbreak{}Credit.\allowbreak{}prior\allowbreak{}Risk}
\end{itemize}
\clearpage
\hypertarget{paper-prerequisite-section-5ef5ef0364b6939c}{}
\section*{Paper-specific semantic prerequisites}
\hypertarget{paper-prerequisite-6956aaccf5fffc8f}{}
\subsection*{\small\ttfamily\raggedright Galarza26\allowbreak{}Relative\allowbreak{}Credit.\allowbreak{}premium}
\noindent\textbf{Paper-local declaration:}\par
{\footnotesize\ttfamily\raggedright Galarza26\allowbreak{}Relative\allowbreak{}Credit.\allowbreak{}premium\par}
\textbf{Source locator:} {\footnotesize\raggedright cited publication, lines 316-\allowbreak{}334\par}
\paragraph{Verbatim paper-source connection}
\begin{ReviewVerbatim}
Supuesto 3 (Remuneración y problema condicionado). El premio normalizado por desempeño es
R(D) = R̄ + D > 0 y F es un costo por unidad independiente de e. El banco resuelve
máx
e∈[0,u(D,κ)]

π(e; D, κ) = [1 − q 0 (D, κ) + g(κ)e]R(D) − F −

c(κ) 2
e .
2

(2)

El precio, la referencia de pares y κ se mantienen fijos al evaluar desviaciones de esfuerzo.
El costo F no afecta el óptimo de esfuerzo. Un volumen estrictamente positivo que multiplique
todo el objetivo tampoco lo afecta si no depende de e. La participación o decisión de no prestar
no está resuelta: el resultado se refiere a una unidad que el banco ya contempla originar. La remuneración supone que evitar el evento adverso mejora el pago esperado. Recuperaciones, seguros y
garantías que varían con precio o esfuerzo exigirían modificar este premio. El horizonte del evento
y el flujo contractual deben ser compatibles en una implementación empírica.
\end{ReviewVerbatim}



\paragraph{Lean-expanded paper semantic target}
\begin{ReviewVerbatim}
type:
ℝ → ℝ → ℝ

value:
fun (Rbar D : ℝ) => Rbar + D
\end{ReviewVerbatim}


\paragraph{Recorded source-to-declaration screening}
\textbf{Verdict:} \texttt{matches}
\\
{\raggedright
\textbf{Reason:} The selected S3 formula is R(D)=R̄+D, and the expanded target is exactly fun (Rbar D) ↦ Rbar+D.\allowbreak{} The displayed conditionedPayoff prerequisite supplies the source's R(D)>0 restriction over the admissible interval, and every displayed regime, optimizer, and proposition use receives the same premium expression;\allowbreak{} no formula, domain case, or conclusion is changed.\allowbreak{}
\par}
\reviewmatch{Matches source input}
\noindent\textbf{Reviewer annotation}\par
\reviewerbox
\clearpage
\hypertarget{paper-prerequisite-e8657fe6a37f000a}{}
\subsection*{\small\ttfamily\raggedright Galarza26\allowbreak{}Relative\allowbreak{}Credit.\allowbreak{}prior\allowbreak{}Risk}
\noindent\textbf{Paper-local declaration:}\par
{\footnotesize\ttfamily\raggedright Galarza26\allowbreak{}Relative\allowbreak{}Credit.\allowbreak{}prior\allowbreak{}Risk\par}
\textbf{Source locator:} {\footnotesize\raggedright cited publication, lines 283-\allowbreak{}284\par}
\paragraph{Verbatim paper-source connection}
\begin{ReviewVerbatim}
Supuesto 1 (Selección y dominio). Para cada κ fijo, a = a(κ) > 0, q̄ ∈ (0, 1) y D pertenece a un
intervalo compacto D. La probabilidad previa es q 0 (D, κ) = q̄ + a(κ)D ∈ [0, 1] en todo D.
\end{ReviewVerbatim}



\paragraph{Lean-expanded paper semantic target}
\begin{ReviewVerbatim}
type:
ℝ → ℝ → ℝ → ℝ

value:
fun (qbar a D : ℝ) => qbar + a * D
\end{ReviewVerbatim}


\paragraph{Recorded source-to-declaration screening}
\textbf{Verdict:} \texttt{matches}
\\
{\raggedright
\textbf{Reason:} Compared Assumption 1's prior probability q0(D,kappa) = qbar + a(kappa)D with the expanded target qbar + a*D for fixed kappa.\allowbreak{} The displayed selectionAndDomain and proposition uses impose qbar in (0,1), a>0, D in the compact interval, and the resulting value in [0,1], so the general real-\allowbreak{}valued definition is used on precisely the source domain.\allowbreak{}
\par}
\reviewmatch{Matches source input}
\noindent\textbf{Reviewer annotation}\par
\reviewerbox
\clearpage
\hypertarget{paper-prerequisite-3f518444e1a02f7c}{}
\subsection*{\small\ttfamily\raggedright Galarza26\allowbreak{}Relative\allowbreak{}Credit.\allowbreak{}screening\allowbreak{}Cost}
\noindent\textbf{Paper-local declaration:}\par
{\footnotesize\ttfamily\raggedright Galarza26\allowbreak{}Relative\allowbreak{}Credit.\allowbreak{}screening\allowbreak{}Cost\par}
\textbf{Source locator:} {\footnotesize\raggedright cited publication, lines 299-\allowbreak{}302\par}
\paragraph{Verbatim paper-source connection}
\begin{ReviewVerbatim}
Supuesto 2 (Tecnología, capacidad y costo). g = g(κ) > 0, c = c(κ) > 0 y ē ≥ 0 son finitos.
Para un volumen normalizado de créditos, q(D, e, κ) = q 0 (D, κ) − g(κ)e y el costo es c(κ)e2 /2. La
elección factible cumple
0 ≤ e ≤ u(D, κ) := mı́n{ē, q 0 (D, κ)/g(κ)}.
\end{ReviewVerbatim}



\paragraph{Lean-expanded paper semantic target}
\begin{ReviewVerbatim}
type:
ℝ → ℝ → ℝ

value:
fun (c e : ℝ) => c * e ^ 2 / 2
\end{ReviewVerbatim}


\paragraph{Recorded source-to-declaration screening}
\textbf{Verdict:} \texttt{matches}
\\
{\raggedright
\textbf{Reason:} Compared Assumption 2's normalized quadratic cost c(kappa)e\textasciicircum{}2/\allowbreak{}2 with the expanded target c*e\textasciicircum{}2/\allowbreak{}2;\allowbreak{} the coefficient, exponent, divisor, and effort binder agree exactly.\allowbreak{} Its bankProfit and continuation-\allowbreak{}problem uses hold c fixed during effort comparison and technologyCapacityCost supplies c>0.\allowbreak{}
\par}
\reviewmatch{Matches source input}
\noindent\textbf{Reviewer annotation}\par
\reviewerbox
\clearpage
\hypertarget{paper-prerequisite-646bdd87bd4c2b85}{}
\subsection*{\small\ttfamily\raggedright Galarza26\allowbreak{}Relative\allowbreak{}Credit.\allowbreak{}bank\allowbreak{}Profit}
\noindent\textbf{Paper-local declaration:}\par
{\footnotesize\ttfamily\raggedright Galarza26\allowbreak{}Relative\allowbreak{}Credit.\allowbreak{}bank\allowbreak{}Profit\par}
\textbf{Source locator:} {\footnotesize\raggedright cited publication, lines 316-\allowbreak{}334\par}
\paragraph{Verbatim paper-source connection}
\begin{ReviewVerbatim}
Supuesto 3 (Remuneración y problema condicionado). El premio normalizado por desempeño es
R(D) = R̄ + D > 0 y F es un costo por unidad independiente de e. El banco resuelve
máx
e∈[0,u(D,κ)]

π(e; D, κ) = [1 − q 0 (D, κ) + g(κ)e]R(D) − F −

c(κ) 2
e .
2

(2)

El precio, la referencia de pares y κ se mantienen fijos al evaluar desviaciones de esfuerzo.
El costo F no afecta el óptimo de esfuerzo. Un volumen estrictamente positivo que multiplique
todo el objetivo tampoco lo afecta si no depende de e. La participación o decisión de no prestar
no está resuelta: el resultado se refiere a una unidad que el banco ya contempla originar. La remuneración supone que evitar el evento adverso mejora el pago esperado. Recuperaciones, seguros y
garantías que varían con precio o esfuerzo exigirían modificar este premio. El horizonte del evento
y el flujo contractual deben ser compatibles en una implementación empírica.
\end{ReviewVerbatim}



\paragraph{Lean-expanded paper semantic target}
\begin{ReviewVerbatim}
type:
ℝ → ℝ → ℝ → ℝ → ℝ → ℝ → ℝ → ℝ → ℝ

value:
fun (qbar a g c Rbar F D e : ℝ) =>
  (1 - Galarza26RelativeCredit.priorRisk qbar a D + g * e) * Galarza26RelativeCredit.premium Rbar D - F -
    Galarza26RelativeCredit.screeningCost c e
\end{ReviewVerbatim}

\paragraph{Direct retained dependencies}
\begin{itemize}\raggedright
\item \hyperlink{paper-prerequisite-6956aaccf5fffc8f}{Galarza26\allowbreak{}Relative\allowbreak{}Credit.\allowbreak{}premium}
\item \hyperlink{paper-prerequisite-e8657fe6a37f000a}{Galarza26\allowbreak{}Relative\allowbreak{}Credit.\allowbreak{}prior\allowbreak{}Risk}
\item \hyperlink{paper-prerequisite-3f518444e1a02f7c}{Galarza26\allowbreak{}Relative\allowbreak{}Credit.\allowbreak{}screening\allowbreak{}Cost}
\end{itemize}
\paragraph{Recorded source-to-declaration screening}
\textbf{Verdict:} \texttt{matches}
\\
{\raggedright
\textbf{Reason:} The selected S3 span gives π(e;\allowbreak{}D,κ) = [1-\allowbreak{}q₀(D,κ)+g(κ)e](R̄+D)-\allowbreak{}F-\allowbreak{}c(κ)e²/\allowbreak{}2 with D and κ fixed during an effort comparison.\allowbreak{} The expanded target substitutes priorRisk = qbar+aD, premium = Rbar+D, and screeningCost = c e²/\allowbreak{}2, yielding exactly that payoff for fixed scalar primitives;\allowbreak{} its displayed paper use is through the conditioned problem under the source-\allowbreak{}domain assumptions, so the unrestricted raw function introduces no different source conclusion.\allowbreak{}
\par}
\reviewmatch{Matches source input}
\noindent\textbf{Reviewer annotation}\par
\reviewerbox
\clearpage
\hypertarget{paper-prerequisite-b0571b2c4a78eab5}{}
\subsection*{\small\ttfamily\raggedright Galarza26\allowbreak{}Relative\allowbreak{}Credit.\allowbreak{}cohort\allowbreak{}Weight}
\noindent\textbf{Paper-local declaration:}\par
{\footnotesize\ttfamily\raggedright Galarza26\allowbreak{}Relative\allowbreak{}Credit.\allowbreak{}cohort\allowbreak{}Weight\par}
\textbf{Source locator:} {\footnotesize\raggedright cited publication, lines 523-\allowbreak{}525\par}
\paragraph{Verbatim paper-source connection}
\begin{ReviewVerbatim}
Fijemos un horizonte h. B0 > 0 es el saldo de créditos antiguos todavía vigente en esa fecha y
m0 ∈ [0, 1] su proporción en el estado de atraso elegido. B(D) > 0 es el saldo de la cohorte nueva
que permanece en cartera a h, después de amortización normal. Denote w(D) = B(D)/(B0 +B(D)).
\end{ReviewVerbatim}



\paragraph{Lean-expanded paper semantic target}
\begin{ReviewVerbatim}
type:
ℝ → (ℝ → ℝ) → ℝ → ℝ

value:
fun (B0 : ℝ) (B : ℝ → ℝ) (D : ℝ) => B D / (B0 + B D)
\end{ReviewVerbatim}


\paragraph{Recorded source-to-declaration screening}
\textbf{Verdict:} \texttt{matches}
\\
{\raggedright
\textbf{Reason:} Compared the stock-\allowbreak{}model definition w(D) = B(D)/\allowbreak{}(B0+B(D)) with the expanded target B D/\allowbreak{}(B0+B D);\allowbreak{} the binders and formula coincide.\allowbreak{} In the displayed proposition-\allowbreak{}3 use, stockAccountingConditions supplies B0 > 0 and B(D) > 0 on the paper interval, so the more generally typed definition specializes to the source-\allowbreak{}admissible nonzero denominator.\allowbreak{}
\par}
\reviewmatch{Matches source input}
\noindent\textbf{Reviewer annotation}\par
\reviewerbox
\clearpage
\hypertarget{paper-prerequisite-3564270d1ab44650}{}
\subsection*{\small\ttfamily\raggedright Galarza26\allowbreak{}Relative\allowbreak{}Credit.\allowbreak{}feasible\allowbreak{}Efforts}
\noindent\textbf{Paper-local declaration:}\par
{\footnotesize\ttfamily\raggedright Galarza26\allowbreak{}Relative\allowbreak{}Credit.\allowbreak{}feasible\allowbreak{}Efforts\par}
\textbf{Source locator:} {\footnotesize\raggedright cited publication, lines 299-\allowbreak{}302\par}
\paragraph{Verbatim paper-source connection}
\begin{ReviewVerbatim}
Supuesto 2 (Tecnología, capacidad y costo). g = g(κ) > 0, c = c(κ) > 0 y ē ≥ 0 son finitos.
Para un volumen normalizado de créditos, q(D, e, κ) = q 0 (D, κ) − g(κ)e y el costo es c(κ)e2 /2. La
elección factible cumple
0 ≤ e ≤ u(D, κ) := mı́n{ē, q 0 (D, κ)/g(κ)}.
\end{ReviewVerbatim}



\paragraph{Lean-expanded paper semantic target}
\begin{ReviewVerbatim}
type:
ℝ → ℝ → ℝ → ℝ → ℝ → ℝ → Prop

value:
fun (qbar a g E D a_1 : ℝ) => Set.Icc 0 (min E (Galarza26RelativeCredit.priorRisk qbar a D / g)) a_1
\end{ReviewVerbatim}

\paragraph{Direct retained dependencies}
\begin{itemize}\raggedright
\item \hyperlink{paper-prerequisite-e8657fe6a37f000a}{Galarza26\allowbreak{}Relative\allowbreak{}Credit.\allowbreak{}prior\allowbreak{}Risk}
\end{itemize}
\paragraph{Recorded source-to-declaration screening}
\textbf{Verdict:} \texttt{matches}
\\
{\raggedright
\textbf{Reason:} Compared Assumption 2's 0 <= e <= min(E,q0(D,kappa)/\allowbreak{}g(kappa)) with the expanded membership target Set.\allowbreak{}Icc 0 (min E (priorRisk qbar a D/\allowbreak{}g)).\allowbreak{} With priorRisk expanded to qbar+a*D, the bounds, divisions, minimum, and closed endpoints coincide, and proposition 1 uses the set under g>0, E>=0, and q0 in [0,1].\allowbreak{}
\par}
\reviewmatch{Matches source input}
\noindent\textbf{Reviewer annotation}\par
\reviewerbox
\clearpage
\hypertarget{paper-prerequisite-a6fae381e2d9a684}{}
\subsection*{\small\ttfamily\raggedright Galarza26\allowbreak{}Relative\allowbreak{}Credit.\allowbreak{}conditioned\allowbreak{}Continuation\allowbreak{}Problem}
\noindent\textbf{Paper-local declaration:}\par
{\footnotesize\ttfamily\raggedright Galarza26\allowbreak{}Relative\allowbreak{}Credit.\allowbreak{}conditioned\allowbreak{}Continuation\allowbreak{}Problem\par}
\textbf{Source locator:} {\footnotesize\raggedright cited publication, lines 211-\allowbreak{}233\par}
\paragraph{Verbatim paper-source connection}
\begin{ReviewVerbatim}
El único agente que optimiza en este modelo es un banco neutral al riesgo. Los solicitantes
y los bancos pares determinan una calidad y una referencia reducidas; no se modela su juego
estratégico. El supervisor observa D y más tarde un ratio de cartera, pero no el esfuerzo ni toda la
composición de riesgo. La opacidad κ ∈ [0, 1] es una característica exógena del segmento: un valor
mayor representa menor información verificable inicial. No se estima este índice ni se equipara
automáticamente microempresa con opacidad alta.
La cronología es: primero se fijan segmento, tecnología y precio observado; después se forma el
conjunto de solicitantes y el banco observa su riesgo medio; luego elige esfuerzo antes de originar;
finalmente se realiza el estado crediticio y se observa la cartera al horizonte h. El banco conoce los
parámetros al decidir. La incertidumbre de desempeño permanece y se integra en su objetivo. El
precio es un argumento del problema de continuación, no una variable resuelta en equilibrio.
q 0 y q son probabilidades sin unidades, referidas al mismo evento y horizonte para una cohorte
estandarizada. e y κ son índices. D se mide en unidades de tasa anual decimal. Los pagos y costos
se normalizan por sol de crédito. La sensibilidad monetaria del premio de desempeño a una unidad
4


[form-feed in source extraction]
de D se normaliza a uno. Esto es una simplificación de remuneración, no una conversión exacta de
tasas efectivas a cualquier plazo contractual.
Conviene separar tres operaciones que aparecerán en las pruebas. Al elegir e, el banco mantiene
fijos D y κ: compara esfuerzos para un mismo precio y segmento. Al calcular ∂D q ∗ , compara
problemas con precios cercanos, dejando fijo el segmento y permitiendo que cambie el esfuerzo
óptimo. Solo en la Proposición 2 se comparan distintos valores de κ, siempre al mismo precio D0 .
La prima de M ′ se refiere a D; las primas de g ′ y c′ en el análisis de opacidad se refieren a κ.
\end{ReviewVerbatim}



\paragraph{Lean-expanded paper semantic target}
\begin{ReviewVerbatim}
type:
ℝ → ℝ → ℝ → ℝ → ℝ → ℝ → ℝ → ℝ → Galarza26RelativeCredit.ContinuationProblem

value:
fun (qbar a g c Rbar F E D : ℝ) =>
  { feasible := Galarza26RelativeCredit.feasibleEfforts qbar a g E D,
    payoff := fun (e : ℝ) => Galarza26RelativeCredit.bankProfit qbar a g c Rbar F D e }
\end{ReviewVerbatim}

\paragraph{Direct retained dependencies}
\begin{itemize}\raggedright
\item \hyperlink{governing-declaration-34d15e1983b4c930}{Galarza26\allowbreak{}Relative\allowbreak{}Credit.\allowbreak{}Continuation\allowbreak{}Problem}
\item \hyperlink{paper-prerequisite-646bdd87bd4c2b85}{Galarza26\allowbreak{}Relative\allowbreak{}Credit.\allowbreak{}bank\allowbreak{}Profit}
\item \hyperlink{paper-prerequisite-3564270d1ab44650}{Galarza26\allowbreak{}Relative\allowbreak{}Credit.\allowbreak{}feasible\allowbreak{}Efforts}
\end{itemize}
\paragraph{Recorded source-to-declaration screening}
\textbf{Verdict:} \texttt{matches}
\\
{\raggedright
\textbf{Reason:} The governing span makes the risk-\allowbreak{}neutral bank the sole optimizer, fixes segment, technology, observed price D, and κ before effort is varied, and expressly treats price as an argument of a continuation problem.\allowbreak{} The expanded target is exactly a one-\allowbreak{}shot feasible-\allowbreak{}set/\allowbreak{}payoff record:\allowbreak{} feasibleEfforts expands to [0,min E ((qbar+aD)/\allowbreak{}g)] and the payoff is e ↦ bankProfit at fixed D and fixed scalar primitives.\allowbreak{} Its Proposition 1 use ranges D over the admissible interval under the displayed selection/\allowbreak{}domain, technology/\allowbreak{}capacity/\allowbreak{}cost, and positive-\allowbreak{}premium premises, preserving the source chronology and conditioned optimization domain.\allowbreak{}
\par}
\reviewmatch{Matches source input}
\noindent\textbf{Reviewer annotation}\par
\reviewerbox
\clearpage
\hypertarget{paper-prerequisite-0bf7e0c122a51179}{}
\subsection*{\small\ttfamily\raggedright Galarza26\allowbreak{}Relative\allowbreak{}Credit.\allowbreak{}conditioned\allowbreak{}Payoff}
\noindent\textbf{Paper-local declaration:}\par
{\footnotesize\ttfamily\raggedright Galarza26\allowbreak{}Relative\allowbreak{}Credit.\allowbreak{}conditioned\allowbreak{}Payoff\par}
\textbf{Source locator:} {\footnotesize\raggedright cited publication, lines 316-\allowbreak{}334\par}
\paragraph{Verbatim paper-source connection}
\begin{ReviewVerbatim}
Supuesto 3 (Remuneración y problema condicionado). El premio normalizado por desempeño es
R(D) = R̄ + D > 0 y F es un costo por unidad independiente de e. El banco resuelve
máx
e∈[0,u(D,κ)]

π(e; D, κ) = [1 − q 0 (D, κ) + g(κ)e]R(D) − F −

c(κ) 2
e .
2

(2)

El precio, la referencia de pares y κ se mantienen fijos al evaluar desviaciones de esfuerzo.
El costo F no afecta el óptimo de esfuerzo. Un volumen estrictamente positivo que multiplique
todo el objetivo tampoco lo afecta si no depende de e. La participación o decisión de no prestar
no está resuelta: el resultado se refiere a una unidad que el banco ya contempla originar. La remuneración supone que evitar el evento adverso mejora el pago esperado. Recuperaciones, seguros y
garantías que varían con precio o esfuerzo exigirían modificar este premio. El horizonte del evento
y el flujo contractual deben ser compatibles en una implementación empírica.
\end{ReviewVerbatim}



\paragraph{Lean-expanded paper semantic target}
\begin{ReviewVerbatim}
type:
ℝ → ℝ → ℝ → Prop

value:
fun (Rbar lower upper : ℝ) => ∀ D ∈ Set.Icc lower upper, 0 < Galarza26RelativeCredit.premium Rbar D
\end{ReviewVerbatim}

\paragraph{Direct retained dependencies}
\begin{itemize}\raggedright
\item \hyperlink{paper-prerequisite-6956aaccf5fffc8f}{Galarza26\allowbreak{}Relative\allowbreak{}Credit.\allowbreak{}premium}
\end{itemize}
\paragraph{Recorded source-to-declaration screening}
\textbf{Verdict:} \texttt{matches}
\\
{\raggedright
\textbf{Reason:} The selected S3 span states R(D)=R̄+D>0.\allowbreak{} The expanded target requires 0 < premium Rbar D for every D in the admissible closed interval, and premium expands to Rbar+D;\allowbreak{} all three displayed proposition uses assume this same interval-\allowbreak{}wide condition.\allowbreak{} Thus the binder, domain restriction, formula, and strict positivity premise match the source.\allowbreak{}
\par}
\reviewmatch{Matches source input}
\noindent\textbf{Reviewer annotation}\par
\reviewerbox
\clearpage
\hypertarget{paper-prerequisite-edcf6b4879cafd1b}{}
\subsection*{\small\ttfamily\raggedright Galarza26\allowbreak{}Relative\allowbreak{}Credit.\allowbreak{}screening\allowbreak{}Strength}
\noindent\textbf{Paper-local declaration:}\par
{\footnotesize\ttfamily\raggedright Galarza26\allowbreak{}Relative\allowbreak{}Credit.\allowbreak{}screening\allowbreak{}Strength\par}
\textbf{Source locator:} {\footnotesize\raggedright cited publication, lines 479\par}
\paragraph{Verbatim paper-source connection}
\begin{ReviewVerbatim}
Defina H(κ) = g(κ)2 /c(κ) y β(κ) = a(κ) − H(κ).
\end{ReviewVerbatim}



\paragraph{Lean-expanded paper semantic target}
\begin{ReviewVerbatim}
type:
ℝ → ℝ → ℝ

value:
fun (g c : ℝ) => g ^ 2 / c
\end{ReviewVerbatim}


\paragraph{Recorded source-to-declaration screening}
\textbf{Verdict:} \texttt{matches}
\\
{\raggedright
\textbf{Reason:} Compared H(kappa) = g(kappa)\textasciicircum{}2/\allowbreak{}c(kappa) with the expanded target g\textasciicircum{}2/\allowbreak{}c.\allowbreak{} The displayed crossing and proposition-\allowbreak{}2 contexts evaluate it pointwise with c(kappa)>0 from the source technology assumptions, and screeningAdvantage then forms exactly a-\allowbreak{}H.\allowbreak{}
\par}
\reviewmatch{Matches source input}
\noindent\textbf{Reviewer annotation}\par
\reviewerbox
\clearpage
\hypertarget{paper-prerequisite-29fcb63a4a0578b3}{}
\subsection*{\small\ttfamily\raggedright Galarza26\allowbreak{}Relative\allowbreak{}Credit.\allowbreak{}screening\allowbreak{}Advantage}
\noindent\textbf{Paper-local declaration:}\par
{\footnotesize\ttfamily\raggedright Galarza26\allowbreak{}Relative\allowbreak{}Credit.\allowbreak{}screening\allowbreak{}Advantage\par}
\textbf{Source locator:} {\footnotesize\raggedright cited publication, lines 479\par}
\paragraph{Verbatim paper-source connection}
\begin{ReviewVerbatim}
Defina H(κ) = g(κ)2 /c(κ) y β(κ) = a(κ) − H(κ).
\end{ReviewVerbatim}



\paragraph{Lean-expanded paper semantic target}
\begin{ReviewVerbatim}
type:
ℝ → ℝ → ℝ → ℝ

value:
fun (a g c : ℝ) => a - Galarza26RelativeCredit.screeningStrength g c
\end{ReviewVerbatim}

\paragraph{Direct retained dependencies}
\begin{itemize}\raggedright
\item \hyperlink{paper-prerequisite-edcf6b4879cafd1b}{Galarza26\allowbreak{}Relative\allowbreak{}Credit.\allowbreak{}screening\allowbreak{}Strength}
\end{itemize}
\paragraph{Recorded source-to-declaration screening}
\textbf{Verdict:} \texttt{matches}
\\
{\raggedright
\textbf{Reason:} Compared beta(kappa) = a(kappa) -\allowbreak{} H(kappa) with the expanded target a -\allowbreak{} screeningStrength g c, and the displayed screeningStrength expands to g\textasciicircum{}2/\allowbreak{}c.\allowbreak{} CrossingConditions and proposition 2 apply this exact quantity pointwise on kappa in [0,1] for endpoint signs, strict decrease, and derivative conclusions.\allowbreak{}
\par}
\reviewmatch{Matches source input}
\noindent\textbf{Reviewer annotation}\par
\reviewerbox
\clearpage
\hypertarget{paper-prerequisite-4f59b48de0be00ac}{}
\subsection*{\small\ttfamily\raggedright Galarza26\allowbreak{}Relative\allowbreak{}Credit.\allowbreak{}crossing\allowbreak{}Conditions}
\noindent\textbf{Paper-local declaration:}\par
{\footnotesize\ttfamily\raggedright Galarza26\allowbreak{}Relative\allowbreak{}Credit.\allowbreak{}crossing\allowbreak{}Conditions\par}
\textbf{Source locator:} {\footnotesize\raggedright cited publication, lines 480-\allowbreak{}483\par}
\paragraph{Verbatim paper-source connection}
\begin{ReviewVerbatim}
Supuesto 4 (Condiciones suficientes para el cruce). En κ ∈ [0, 1], a, g, c son continuas y positivas.
La función a − H es estrictamente decreciente y satisface a(0) > H(0) y a(1) < H(1). Al precio D0
que se compara se cumplen estrictamente las dos condiciones de interioridad de la primera línea de
(4) para todo κ.
\end{ReviewVerbatim}



\paragraph{Lean-expanded paper semantic target}
\begin{ReviewVerbatim}
type:
ℝ → ℝ → ℝ → ℝ → (ℝ → ℝ) → (ℝ → ℝ) → (ℝ → ℝ) → Prop

value:
fun (qbar Rbar E D0 : ℝ) (a g c : ℝ → ℝ) =>
  ContinuousOn a (Set.Icc 0 1) ∧
    ContinuousOn g (Set.Icc 0 1) ∧
      ContinuousOn c (Set.Icc 0 1) ∧
        (∀ k ∈ Set.Icc 0 1, 0 < a k ∧ 0 < g k ∧ 0 < c k) ∧
          StrictAntiOn (fun (k : ℝ) => Galarza26RelativeCredit.screeningAdvantage (a k) (g k) (c k)) (Set.Icc 0 1) ∧
            0 < Galarza26RelativeCredit.screeningAdvantage (a 0) (g 0) (c 0) ∧
              Galarza26RelativeCredit.screeningAdvantage (a 1) (g 1) (c 1) < 0 ∧
                ∀ k ∈ Set.Icc 0 1, Galarza26RelativeCredit.interiorRegime qbar (a k) (g k) (c k) Rbar E D0
\end{ReviewVerbatim}

\paragraph{Direct retained dependencies}
\begin{itemize}\raggedright
\item \hyperlink{governing-declaration-5f246d7c1bbfcd61}{Galarza26\allowbreak{}Relative\allowbreak{}Credit.\allowbreak{}interior\allowbreak{}Regime}
\item \hyperlink{paper-prerequisite-29fcb63a4a0578b3}{Galarza26\allowbreak{}Relative\allowbreak{}Credit.\allowbreak{}screening\allowbreak{}Advantage}
\end{itemize}
\paragraph{Recorded source-to-declaration screening}
\textbf{Verdict:} \texttt{matches}
\\
{\raggedright
\textbf{Reason:} Compared every clause of Assumption 4 with the expanded target on kappa in [0,1]:\allowbreak{} continuity and positivity of a,g,c, strict decrease of a-\allowbreak{}g\textasciicircum{}2/\allowbreak{}c, positive and negative endpoint signs, and both strict interior inequalities at D0 for every kappa are all present.\allowbreak{} The displayed proposition-\allowbreak{}2 use also places D0 in the price interval and supplies S1-\allowbreak{}S3 for every kappa.\allowbreak{}
\par}
\reviewmatch{Matches source input}
\noindent\textbf{Reviewer annotation}\par
\reviewerbox
\clearpage
\hypertarget{paper-prerequisite-527f5e3e1ca86ece}{}
\subsection*{\small\ttfamily\raggedright Galarza26\allowbreak{}Relative\allowbreak{}Credit.\allowbreak{}impaired\allowbreak{}Balance}
\noindent\textbf{Paper-local declaration:}\par
{\footnotesize\ttfamily\raggedright Galarza26\allowbreak{}Relative\allowbreak{}Credit.\allowbreak{}impaired\allowbreak{}Balance\par}
\textbf{Source locator:} {\footnotesize\raggedright cited publication, lines 530-\allowbreak{}535\par}
\paragraph{Verbatim paper-source connection}
\begin{ReviewVerbatim}
Supuesto 5 (Contabilidad y maduración). Al horizonte, el saldo de cada cohorte es determinista
o se trata como exposición de una población continua. El riesgo q ∗ (D) coincide con la fracción
esperada de su saldo en el evento crediticio elegido, condicional a seguir en cartera. La amortización
y el tamaño del crédito no seleccionan riesgo adicional no recogido en q ∗ . B0 , m0 son constantes
al variar D; B(D) > 0 es diferenciable. No hay castigos, ventas de cartera ni reclasificaciones que
alteren esta identidad en el ejercicio.

[Next verbatim source excerpt]

Fijemos un horizonte h. B0 > 0 es el saldo de créditos antiguos todavía vigente en esa fecha y
m0 ∈ [0, 1] su proporción en el estado de atraso elegido. B(D) > 0 es el saldo de la cohorte nueva
que permanece en cartera a h, después de amortización normal. Denote w(D) = B(D)/(B0 +B(D)).
\end{ReviewVerbatim}



\paragraph{Lean-expanded paper semantic target}
\begin{ReviewVerbatim}
type:
ℝ → ℝ → (ℝ → ℝ) → (ℝ → ℝ) → ℝ → ℝ

value:
fun (B0 m0 : ℝ) (B q : ℝ → ℝ) (D : ℝ) => B0 * m0 + B D * q D
\end{ReviewVerbatim}


\paragraph{Recorded source-to-declaration screening}
\textbf{Verdict:} \texttt{matches}
\\
{\raggedright
\textbf{Reason:} The selected source makes m0 the old cohort's event proportion and q*(D) the expected event fraction of the retained new-\allowbreak{}cohort balance, while excluding write-\allowbreak{}offs, sales, reclassifications, and additional risk-\allowbreak{}selection terms.\allowbreak{} The Lean value B0 * m0 + B D * q D is therefore exactly the sum of the two impaired cohort balances;\allowbreak{} the displayed portfolio ratio divides that sum by B0 + B D, and proposition3Spec specializes q to optimalRisk and uses the same weighted two-\allowbreak{}cohort identity on the admissible interval.\allowbreak{} Allowing arbitrary real arguments in the reusable definition does not remove any source-\allowbreak{}admissible case.\allowbreak{}
\par}
\reviewmatch{Matches source input}
\noindent\textbf{Reviewer annotation}\par
\reviewerbox
\clearpage
\hypertarget{paper-prerequisite-3b6383e0216b3c72}{}
\subsection*{\small\ttfamily\raggedright Galarza26\allowbreak{}Relative\allowbreak{}Credit.\allowbreak{}relative\allowbreak{}Price}
\noindent\textbf{Paper-local declaration:}\par
{\footnotesize\ttfamily\raggedright Galarza26\allowbreak{}Relative\allowbreak{}Credit.\allowbreak{}relative\allowbreak{}Price\par}
\textbf{Source locator:} {\footnotesize\raggedright cited publication, lines 127-\allowbreak{}156\par}
\paragraph{Verbatim paper-source connection}
\begin{ReviewVerbatim}
Sea j un banco, s una celda de tipo de crédito, modalidad o tramo de plazo y moneda, y t la
fecha del reporte. Denote por ijst su tasa anual en decimales y por bst una base común en la misma
moneda. Sea P−j,st el conjunto de otros bancos con observación válida y n−j,st su tamaño. Defino
Djst = (ijst − bst ) −

1
n−j,st

X

(ikst − bst ) = ijst −

k∈P−j,st

1

X

n−j,st

k

ikst .

(1)

El banco evaluado no entra en su referencia. D = 0.02 significa dos puntos porcentuales por encima
del promedio de pares. Se usa el promedio simple, coherente con el plan y transparente cuando no
se observan desembolsos por tramo. No es la tasa promedio del sistema ponderada por operaciones.
Se exigirían al menos tres pares para la aplicación empírica.
\end{ReviewVerbatim}



\paragraph{Lean-expanded paper semantic target}
\begin{ReviewVerbatim}
type:
{n : ℕ} → ℝ → ℝ → (Fin (n + 3) → ℝ) → ℝ

value:
fun {n : ℕ} (ownRate baseRate : ℝ) (peerRates : Fin (n + 3) → ℝ) =>
  ownRate - baseRate - (∑ k : Fin (n + 3), (peerRates k - baseRate)) / ↑(Fintype.card (Fin (n + 3)))
\end{ReviewVerbatim}


\paragraph{Recorded source-to-declaration screening}
\textbf{Verdict:} \texttt{matches}
\\
{\raggedright
\textbf{Reason:} The selected equation (1) defines the own-\allowbreak{}minus-\allowbreak{}base spread minus the simple average of other-\allowbreak{}bank spreads, excludes the evaluated bank, and requires at least three peers.\allowbreak{} The expanded target supplies ownRate separately and indexes peerRates by Fin(n+3), so it enforces at least three peer observations and divides their unweighted sum by exactly their cardinality.\allowbreak{} Its formula and displayed algebraic use therefore match the source definition.\allowbreak{}
\par}
\reviewmatch{Matches source input}
\noindent\textbf{Reviewer annotation}\par
\reviewerbox
\clearpage
\hypertarget{paper-prerequisite-38fc4c8d2be8054c}{}
\subsection*{\small\ttfamily\raggedright Galarza26\allowbreak{}Relative\allowbreak{}Credit.\allowbreak{}relative\allowbreak{}Price\_\allowbreak{}eq\_\allowbreak{}rate\_\allowbreak{}gap}
\noindent\textbf{Paper-local declaration:}\par
{\footnotesize\ttfamily\raggedright Galarza26\allowbreak{}Relative\allowbreak{}Credit.\allowbreak{}relative\allowbreak{}Price\_\allowbreak{}eq\_\allowbreak{}rate\_\allowbreak{}gap\par}
\textbf{Source locator:} {\footnotesize\raggedright cited publication, lines 127-\allowbreak{}156\par}
\paragraph{Verbatim paper-source connection}
\begin{ReviewVerbatim}
Sea j un banco, s una celda de tipo de crédito, modalidad o tramo de plazo y moneda, y t la
fecha del reporte. Denote por ijst su tasa anual en decimales y por bst una base común en la misma
moneda. Sea P−j,st el conjunto de otros bancos con observación válida y n−j,st su tamaño. Defino
Djst = (ijst − bst ) −

1
n−j,st

X

(ikst − bst ) = ijst −

k∈P−j,st

1

X

n−j,st

k

ikst .

(1)

El banco evaluado no entra en su referencia. D = 0.02 significa dos puntos porcentuales por encima
del promedio de pares. Se usa el promedio simple, coherente con el plan y transparente cuando no
se observan desembolsos por tramo. No es la tasa promedio del sistema ponderada por operaciones.
Se exigirían al menos tres pares para la aplicación empírica.
\end{ReviewVerbatim}



\paragraph{Lean-elaborated paper signature}
\begin{ReviewVerbatim}
∀ {n : ℕ} (ownRate baseRate : ℝ) (peerRates : Fin (n + 3) → ℝ),
  Galarza26RelativeCredit.relativePrice ownRate baseRate peerRates =
    ownRate - (∑ k : Fin (n + 3), peerRates k) / ↑(Fintype.card (Fin (n + 3)))
\end{ReviewVerbatim}

\paragraph{Direct retained dependencies}
\begin{itemize}\raggedright
\item \hyperlink{paper-prerequisite-3b6383e0216b3c72}{Galarza26\allowbreak{}Relative\allowbreak{}Credit.\allowbreak{}relative\allowbreak{}Price}
\end{itemize}
\paragraph{Recorded source-to-declaration screening}
\textbf{Verdict:} \texttt{matches}
\\
{\raggedright
\textbf{Reason:} The selected equation (1) states that subtracting the same base from the own rate and every peer rate equals the own rate minus the simple peer-\allowbreak{}rate average.\allowbreak{} The target quantifies over the same relativePrice inputs and concludes precisely that identity;\allowbreak{} Fin(n+3) makes the divisor nonzero, so expanding relativePrice cancels the common base for every source-\allowbreak{}admissible peer set.\allowbreak{}
\par}
\reviewmatch{Matches source input}
\noindent\textbf{Reviewer annotation}\par
\reviewerbox
\clearpage
\hypertarget{paper-prerequisite-e9553b2c455b69dd}{}
\subsection*{\small\ttfamily\raggedright Galarza26\allowbreak{}Relative\allowbreak{}Credit.\allowbreak{}screened\allowbreak{}Risk}
\noindent\textbf{Paper-local declaration:}\par
{\footnotesize\ttfamily\raggedright Galarza26\allowbreak{}Relative\allowbreak{}Credit.\allowbreak{}screened\allowbreak{}Risk\par}
\textbf{Source locator:} {\footnotesize\raggedright cited publication, lines 299-\allowbreak{}302\par}
\paragraph{Verbatim paper-source connection}
\begin{ReviewVerbatim}
Supuesto 2 (Tecnología, capacidad y costo). g = g(κ) > 0, c = c(κ) > 0 y ē ≥ 0 son finitos.
Para un volumen normalizado de créditos, q(D, e, κ) = q 0 (D, κ) − g(κ)e y el costo es c(κ)e2 /2. La
elección factible cumple
0 ≤ e ≤ u(D, κ) := mı́n{ē, q 0 (D, κ)/g(κ)}.
\end{ReviewVerbatim}



\paragraph{Lean-expanded paper semantic target}
\begin{ReviewVerbatim}
type:
ℝ → ℝ → ℝ → ℝ → ℝ → ℝ

value:
fun (qbar a g D e : ℝ) => Galarza26RelativeCredit.priorRisk qbar a D - g * e
\end{ReviewVerbatim}

\paragraph{Direct retained dependencies}
\begin{itemize}\raggedright
\item \hyperlink{paper-prerequisite-e8657fe6a37f000a}{Galarza26\allowbreak{}Relative\allowbreak{}Credit.\allowbreak{}prior\allowbreak{}Risk}
\end{itemize}
\paragraph{Recorded source-to-declaration screening}
\textbf{Verdict:} \texttt{matches}
\\
{\raggedright
\textbf{Reason:} Compared Assumption 2's post-\allowbreak{}screening probability q(D,e,kappa) = q0(D,kappa) -\allowbreak{} g(kappa)e with the expanded target priorRisk qbar a D -\allowbreak{} g*e.\allowbreak{} The displayed feasible-\allowbreak{}set and proposition-\allowbreak{}1 contexts impose the source bounds that keep the optimized risk on the probability domain.\allowbreak{}
\par}
\reviewmatch{Matches source input}
\noindent\textbf{Reviewer annotation}\par
\reviewerbox
\clearpage
\hypertarget{paper-prerequisite-ab0e2f9b84884bc6}{}
\subsection*{\small\ttfamily\raggedright Galarza26\allowbreak{}Relative\allowbreak{}Credit.\allowbreak{}selection\allowbreak{}And\allowbreak{}Domain}
\noindent\textbf{Paper-local declaration:}\par
{\footnotesize\ttfamily\raggedright Galarza26\allowbreak{}Relative\allowbreak{}Credit.\allowbreak{}selection\allowbreak{}And\allowbreak{}Domain\par}
\textbf{Source locator:} {\footnotesize\raggedright cited publication, lines 283-\allowbreak{}284\par}
\paragraph{Verbatim paper-source connection}
\begin{ReviewVerbatim}
Supuesto 1 (Selección y dominio). Para cada κ fijo, a = a(κ) > 0, q̄ ∈ (0, 1) y D pertenece a un
intervalo compacto D. La probabilidad previa es q 0 (D, κ) = q̄ + a(κ)D ∈ [0, 1] en todo D.
\end{ReviewVerbatim}



\paragraph{Lean-expanded paper semantic target}
\begin{ReviewVerbatim}
type:
ℝ → ℝ → ℝ → ℝ → Prop

value:
fun (qbar a lower upper : ℝ) =>
  0 < qbar ∧
    qbar < 1 ∧
      0 < a ∧ lower ≤ upper ∧ ∀ D ∈ Set.Icc lower upper, Galarza26RelativeCredit.priorRisk qbar a D ∈ Set.Icc 0 1
\end{ReviewVerbatim}

\paragraph{Direct retained dependencies}
\begin{itemize}\raggedright
\item \hyperlink{paper-prerequisite-e8657fe6a37f000a}{Galarza26\allowbreak{}Relative\allowbreak{}Credit.\allowbreak{}prior\allowbreak{}Risk}
\end{itemize}
\paragraph{Recorded source-to-declaration screening}
\textbf{Verdict:} \texttt{matches}
\\
{\raggedright
\textbf{Reason:} Compared all of Assumption 1 with the target:\allowbreak{} 0<qbar<1, a>0, a nonempty closed interval lower<=upper, and qbar+a*D in [0,1] for every D in that interval are all explicit.\allowbreak{} Each displayed proposition use invokes this predicate at the relevant fixed kappa, so it specializes exactly to the source's compact-\allowbreak{}domain statement.\allowbreak{}
\par}
\reviewmatch{Matches source input}
\noindent\textbf{Reviewer annotation}\par
\reviewerbox
\clearpage
\hypertarget{paper-prerequisite-de4cba6057d51d83}{}
\subsection*{\small\ttfamily\raggedright Galarza26\allowbreak{}Relative\allowbreak{}Credit.\allowbreak{}stock\allowbreak{}Accounting\allowbreak{}Conditions}
\noindent\textbf{Paper-local declaration:}\par
{\footnotesize\ttfamily\raggedright Galarza26\allowbreak{}Relative\allowbreak{}Credit.\allowbreak{}stock\allowbreak{}Accounting\allowbreak{}Conditions\par}
\textbf{Source locator:} {\footnotesize\raggedright cited publication, lines 530-\allowbreak{}535\par}
\paragraph{Verbatim paper-source connection}
\begin{ReviewVerbatim}
Supuesto 5 (Contabilidad y maduración). Al horizonte, el saldo de cada cohorte es determinista
o se trata como exposición de una población continua. El riesgo q ∗ (D) coincide con la fracción
esperada de su saldo en el evento crediticio elegido, condicional a seguir en cartera. La amortización
y el tamaño del crédito no seleccionan riesgo adicional no recogido en q ∗ . B0 , m0 son constantes
al variar D; B(D) > 0 es diferenciable. No hay castigos, ventas de cartera ni reclasificaciones que
alteren esta identidad en el ejercicio.

[Next verbatim source excerpt]

Fijemos un horizonte h. B0 > 0 es el saldo de créditos antiguos todavía vigente en esa fecha y
m0 ∈ [0, 1] su proporción en el estado de atraso elegido. B(D) > 0 es el saldo de la cohorte nueva
que permanece en cartera a h, después de amortización normal. Denote w(D) = B(D)/(B0 +B(D)).
\end{ReviewVerbatim}



\paragraph{Lean-expanded paper semantic target}
\begin{ReviewVerbatim}
type:
ℝ → ℝ → ℝ → ℝ → (ℝ → ℝ) → Prop

value:
fun (B0 m0 lower upper : ℝ) (B : ℝ → ℝ) =>
  0 < B0 ∧ m0 ∈ Set.Icc 0 1 ∧ ∀ D ∈ Set.Icc lower upper, 0 < B D ∧ DifferentiableAt ℝ B D
\end{ReviewVerbatim}


\paragraph{Recorded source-to-declaration screening}
\textbf{Verdict:} \texttt{matches}
\\
{\raggedright
\textbf{Reason:} For the selected source, B0 and m0 are D-\allowbreak{}independent scalars with B0 > 0 and m0 in [0,1], and B(D) is positive and differentiable as D varies.\allowbreak{} The Lean predicate states precisely those quantitative conditions, pointwise for every D in [lower, upper];\allowbreak{} proposition3Spec assumes the predicate and makes every displayed accounting and derivative use on that same interval.\allowbreak{} The source's deterministic/\allowbreak{}continuous-\allowbreak{}population and no-\allowbreak{}adjustment conventions are reflected in the displayed impaired-\allowbreak{}balance and portfolio-\allowbreak{}ratio formulas rather than being lost from the paper use.\allowbreak{}
\par}
\reviewmatch{Matches source input}
\noindent\textbf{Reviewer annotation}\par
\reviewerbox
\clearpage
\hypertarget{paper-prerequisite-e5f9b7ce782f743a}{}
\subsection*{\small\ttfamily\raggedright Galarza26\allowbreak{}Relative\allowbreak{}Credit.\allowbreak{}technology\allowbreak{}Capacity\allowbreak{}Cost}
\noindent\textbf{Paper-local declaration:}\par
{\footnotesize\ttfamily\raggedright Galarza26\allowbreak{}Relative\allowbreak{}Credit.\allowbreak{}technology\allowbreak{}Capacity\allowbreak{}Cost\par}
\textbf{Source locator:} {\footnotesize\raggedright cited publication, lines 299-\allowbreak{}302\par}
\paragraph{Verbatim paper-source connection}
\begin{ReviewVerbatim}
Supuesto 2 (Tecnología, capacidad y costo). g = g(κ) > 0, c = c(κ) > 0 y ē ≥ 0 son finitos.
Para un volumen normalizado de créditos, q(D, e, κ) = q 0 (D, κ) − g(κ)e y el costo es c(κ)e2 /2. La
elección factible cumple
0 ≤ e ≤ u(D, κ) := mı́n{ē, q 0 (D, κ)/g(κ)}.
\end{ReviewVerbatim}



\paragraph{Lean-expanded paper semantic target}
\begin{ReviewVerbatim}
type:
ℝ → ℝ → ℝ → Prop

value:
fun (g c E : ℝ) => 0 < g ∧ 0 < c ∧ 0 ≤ E
\end{ReviewVerbatim}


\paragraph{Recorded source-to-declaration screening}
\textbf{Verdict:} \texttt{matches}
\\
{\raggedright
\textbf{Reason:} Compared Assumption 2's finite g>0, c>0, and E>=0 with the expanded target containing exactly those inequalities.\allowbreak{} Lean real binders already exclude infinite values, and every displayed proposition use combines this predicate with the source selection and remuneration domains before using the reusable formulas.\allowbreak{}
\par}
\reviewmatch{Matches source input}
\noindent\textbf{Reviewer annotation}\par
\reviewerbox

\clearpage
\hypertarget{source-claim-71d1c57a6546a5a7}{}
\hypertarget{source-section-f10b5f68e4acf3dc}{}
\section*{Source claims}
\subsection*{Review row 1}
\textbf{Source locator:} {\footnotesize\raggedright cited publication, lines 349-\allowbreak{}388\par}
\paragraph{Verbatim source input}
\begin{ReviewVerbatim}
Proposición 1 (Solución global y pendientes por régimen). Bajo los Supuestos 1–3, existe un
único máximo en (2). Es
[U+001A control character]
[U+001B control character]
[U+001A control character]
[U+001A control character]
[U+001B control character][U+001B control character]
gR(D)
q 0 (D)
gR(D)
∗
0
∗
e = mı́n
, ē,
,
q = máx 0, q (D) − g mı́n ē,
.
(3)
c
g
c
Para κ fijo, en cualquier intervalo abierto con desigualdades estrictas que mantengan el régimen,
las pendientes son
[U+F8F1 delimiter glyph]
[U+F8F4 delimiter glyph]a − g 2 /c, gR/c < ē y gR/c < q 0 /g,
∂q ∗ [U+F8F2 delimiter glyph]
= a,
(4)
ē < gR/c y ē < q 0 /g,
[U+F8F4 delimiter glyph]
∂D
[U+F8F3 delimiter glyph]
0
0
0,
q /g < gR/c y q /g < ē.
Si ē = 0, e∗ = 0 y q ∗ = q 0 , con pendiente a en el interior del dominio de D. En empates de
restricciones la solución es continua y se usan derivadas laterales; una derivada ordinaria puede
no existir.

[Next verbatim source excerpt]

primarias consultadas. No es evidencia de que nuestro indicador funcione. Su construcción separa
el deterioro de créditos inicialmente normales durante una ventana de doce meses del ratio contable
de morosidad. Esta diferencia respalda la necesidad de explicitar evento y horizonte; no justifica
equiparar automáticamente nuestro q ∗ con una serie publicada. Clark and West (2007) proporciona
una comparación de modelos predictivos anidados, relevante para un trabajo empírico posterior.
No se calcula aquí una prueba ni un valor p.
Los mecanismos rivales tienen implicaciones que ayudan a diseñar esa prueba. Si domina selección, precios altos anteceden peor calidad de nuevas cosechas, manteniendo constante la evaluación.
Si domina evaluación, el precio alto debe acompañarse de mayor esfuerzo efectivo o mejor documentación y menor riesgo posterior, salvo que la capacidad esté saturada. Si domina composición, el
cambio de signo aparece al agregar al saldo y depende de la participación de colocaciones recientes.
Los costos de fondeo o venta cruzada pueden generar brechas sin estos patrones. Con promedios
públicos será difícil separarlos; harían falta, idealmente, datos de cosechas, rotación, evaluación y
garantías.

4.

Ambiente y decisión del banco

4.1.

Agentes, información, unidades y cronología

El único agente que optimiza en este modelo es un banco neutral al riesgo. Los solicitantes

[Next verbatim source excerpt]

y los bancos pares determinan una calidad y una referencia reducidas; no se modela su juego
estratégico. El supervisor observa D y más tarde un ratio de cartera, pero no el esfuerzo ni toda la

[Next verbatim source excerpt]

conjunto de solicitantes y el banco observa su riesgo medio; luego elige esfuerzo antes de originar;
finalmente se realiza el estado crediticio y se observa la cartera al horizonte h. El banco conoce los
parámetros al decidir. La incertidumbre de desempeño permanece y se integra en su objetivo. El
precio es un argumento del problema de continuación, no una variable resuelta en equilibrio.

[Next verbatim source excerpt]

problemas con precios cercanos, dejando fijo el segmento y permitiendo que cambie el esfuerzo
óptimo. Solo en la Proposición 2 se comparan distintos valores de κ, siempre al mismo precio D0 .
La prima de M ′ se refiere a D; las primas de g ′ y c′ en el análisis de opacidad se refieren a κ.
La tabla siguiente sirve como diccionario. Los objetos de cartera se usan recién en la Proposición 3; no intervienen en la elección inicial del esfuerzo.
Símbolo

Significado

Unidad o restricción

D
κ
q̄, a
q 0 , q, q ∗

[Next verbatim source excerpt]

de morosidad. Esta diferencia respalda la necesidad de explicitar evento y horizonte; no justifica
equiparar automáticamente nuestro q ∗ con una serie publicada. Clark and West (2007) proporciona
una comparación de modelos predictivos anidados, relevante para un trabajo empírico posterior.
No se calcula aquí una prueba ni un valor p.
Los mecanismos rivales tienen implicaciones que ayudan a diseñar esa prueba. Si domina selección, precios altos anteceden peor calidad de nuevas cosechas, manteniendo constante la evaluación.
Si domina evaluación, el precio alto debe acompañarse de mayor esfuerzo efectivo o mejor documentación y menor riesgo posterior, salvo que la capacidad esté saturada. Si domina composición, el
cambio de signo aparece al agregar al saldo y depende de la participación de colocaciones recientes.
Los costos de fondeo o venta cruzada pueden generar brechas sin estos patrones. Con promedios
públicos será difícil separarlos; harían falta, idealmente, datos de cosechas, rotación, evaluación y
garantías.

4.

Ambiente y decisión del banco

4.1.

Agentes, información, unidades y cronología

El único agente que optimiza en este modelo es un banco neutral al riesgo. Los solicitantes
y los bancos pares determinan una calidad y una referencia reducidas; no se modela su juego
estratégico. El supervisor observa D y más tarde un ratio de cartera, pero no el esfuerzo ni toda la

[Next verbatim source excerpt]

composición de riesgo. La opacidad κ ∈ [0, 1] es una característica exógena del segmento: un valor
mayor representa menor información verificable inicial. No se estima este índice ni se equipara

[Next verbatim source excerpt]

parámetros al decidir. La incertidumbre de desempeño permanece y se integra en su objetivo. El
precio es un argumento del problema de continuación, no una variable resuelta en equilibrio.
q 0 y q son probabilidades sin unidades, referidas al mismo evento y horizonte para una cohorte
estandarizada. e y κ son índices. D se mide en unidades de tasa anual decimal. Los pagos y costos

[Next verbatim source excerpt]

La prima de M ′ se refiere a D; las primas de g ′ y c′ en el análisis de opacidad se refieren a κ.
La tabla siguiente sirve como diccionario. Los objetos de cartera se usan recién en la Proposición 3; no intervienen en la elección inicial del esfuerzo.
Símbolo

Significado

Unidad o restricción

D
κ
q̄, a
q 0 , q, q ∗
e, ē, u
g, c

[Next verbatim source excerpt]

El único agente que optimiza en este modelo es un banco neutral al riesgo. Los solicitantes
y los bancos pares determinan una calidad y una referencia reducidas; no se modela su juego
estratégico. El supervisor observa D y más tarde un ratio de cartera, pero no el esfuerzo ni toda la
composición de riesgo. La opacidad κ ∈ [0, 1] es una característica exógena del segmento: un valor
mayor representa menor información verificable inicial. No se estima este índice ni se equipara
automáticamente microempresa con opacidad alta.
La cronología es: primero se fijan segmento, tecnología y precio observado; después se forma el
conjunto de solicitantes y el banco observa su riesgo medio; luego elige esfuerzo antes de originar;
finalmente se realiza el estado crediticio y se observa la cartera al horizonte h. El banco conoce los
parámetros al decidir. La incertidumbre de desempeño permanece y se integra en su objetivo. El
precio es un argumento del problema de continuación, no una variable resuelta en equilibrio.
q 0 y q son probabilidades sin unidades, referidas al mismo evento y horizonte para una cohorte
estandarizada. e y κ son índices. D se mide en unidades de tasa anual decimal. Los pagos y costos
se normalizan por sol de crédito. La sensibilidad monetaria del premio de desempeño a una unidad
4


[form-feed in source extraction]
de D se normaliza a uno. Esto es una simplificación de remuneración, no una conversión exacta de
tasas efectivas a cualquier plazo contractual.
Conviene separar tres operaciones que aparecerán en las pruebas. Al elegir e, el banco mantiene
fijos D y κ: compara esfuerzos para un mismo precio y segmento. Al calcular ∂D q ∗ , compara
problemas con precios cercanos, dejando fijo el segmento y permitiendo que cambie el esfuerzo
óptimo. Solo en la Proposición 2 se comparan distintos valores de κ, siempre al mismo precio D0 .
La prima de M ′ se refiere a D; las primas de g ′ y c′ en el análisis de opacidad se refieren a κ.

[Next verbatim source excerpt]

Supuesto 1 (Selección y dominio). Para cada κ fijo, a = a(κ) > 0, q̄ ∈ (0, 1) y D pertenece a un
intervalo compacto D. La probabilidad previa es q 0 (D, κ) = q̄ + a(κ)D ∈ [0, 1] en todo D.

[Next verbatim source excerpt]

Supuesto 2 (Tecnología, capacidad y costo). g = g(κ) > 0, c = c(κ) > 0 y ē ≥ 0 son finitos.
Para un volumen normalizado de créditos, q(D, e, κ) = q 0 (D, κ) − g(κ)e y el costo es c(κ)e2 /2. La
elección factible cumple
0 ≤ e ≤ u(D, κ) := mı́n{ē, q 0 (D, κ)/g(κ)}.

[Next verbatim source excerpt]

Supuesto 3 (Remuneración y problema condicionado). El premio normalizado por desempeño es
R(D) = R̄ + D > 0 y F es un costo por unidad independiente de e. El banco resuelve
máx
e∈[0,u(D,κ)]

π(e; D, κ) = [1 − q 0 (D, κ) + g(κ)e]R(D) − F −

c(κ) 2
e .
2

(2)

El precio, la referencia de pares y κ se mantienen fijos al evaluar desviaciones de esfuerzo.
El costo F no afecta el óptimo de esfuerzo. Un volumen estrictamente positivo que multiplique
todo el objetivo tampoco lo afecta si no depende de e. La participación o decisión de no prestar
no está resuelta: el resultado se refiere a una unidad que el banco ya contempla originar. La remuneración supone que evitar el evento adverso mejora el pago esperado. Recuperaciones, seguros y
garantías que varían con precio o esfuerzo exigirían modificar este premio. El horizonte del evento
y el flujo contractual deben ser compatibles en una implementación empírica.
\end{ReviewVerbatim}



\paragraph{Expanded PaperInterface specification}
\begin{ReviewVerbatim}
(∀ (qbar a g c Rbar F E lower upper : ℝ),
    Galarza26RelativeCredit.selectionAndDomain qbar a lower upper →
      Galarza26RelativeCredit.technologyCapacityCost g c E →
        Galarza26RelativeCredit.conditionedPayoff Rbar lower upper →
          (∀ D ∈ Set.Icc lower upper,
              ((∃! e : ℝ,
                    (Galarza26RelativeCredit.conditionedContinuationProblem qbar a g c Rbar F E D).IsMaximizer e) ∧
                  (Galarza26RelativeCredit.conditionedContinuationProblem qbar a g c Rbar F E D).IsMaximizer
                    (Galarza26RelativeCredit.optimalEffort qbar a g c Rbar E D)) ∧
                Galarza26RelativeCredit.screenedRisk qbar a g D
                      (Galarza26RelativeCredit.optimalEffort qbar a g c Rbar E D) =
                    Galarza26RelativeCredit.optimalRisk qbar a g c Rbar E D ∧
                  0 ≤ Galarza26RelativeCredit.optimalRisk qbar a g c Rbar E D ∧
                    Galarza26RelativeCredit.optimalRisk qbar a g c Rbar E D ≤ 1) ∧
            (∀ D ∈ Set.Ioo lower upper,
                Galarza26RelativeCredit.interiorRegime qbar a g c Rbar E D →
                  HasDerivAt (Galarza26RelativeCredit.optimalRisk qbar a g c Rbar E) (a - g ^ 2 / c) D) ∧
              (∀ D ∈ Set.Ioo lower upper,
                  Galarza26RelativeCredit.capacityRegime qbar a g c Rbar E D →
                    HasDerivAt (Galarza26RelativeCredit.optimalRisk qbar a g c Rbar E) a D) ∧
                (∀ D ∈ Set.Ioo lower upper,
                    Galarza26RelativeCredit.floorRegime qbar a g c Rbar E D →
                      HasDerivAt (Galarza26RelativeCredit.optimalRisk qbar a g c Rbar E) 0 D) ∧
                  (E = 0 →
                      (∀ D ∈ Set.Icc lower upper,
                          Galarza26RelativeCredit.optimalEffort qbar a g c Rbar E D = 0 ∧
                            Galarza26RelativeCredit.optimalRisk qbar a g c Rbar E D =
                              Galarza26RelativeCredit.priorRisk qbar a D) ∧
                        ∀ D ∈ Set.Ioo lower upper,
                          HasDerivAt (Galarza26RelativeCredit.optimalRisk qbar a g c Rbar E) a D) ∧
                    ContinuousOn (Galarza26RelativeCredit.optimalEffort qbar a g c Rbar E) (Set.Icc lower upper) ∧
                      ContinuousOn (Galarza26RelativeCredit.optimalRisk qbar a g c Rbar E) (Set.Icc lower upper) ∧
                        ∀ D ∈ Set.Icc lower upper,
                          ∃ (left : ℝ) (right : ℝ),
                            HasDerivWithinAt (Galarza26RelativeCredit.optimalRisk qbar a g c Rbar E) left (Set.Iic D)
                                D ∧
                              HasDerivWithinAt (Galarza26RelativeCredit.optimalRisk qbar a g c Rbar E) right (Set.Ici D)
                                D) ∧
  ∃ (qbar' : ℝ) (a' : ℝ) (g' : ℝ) (c' : ℝ) (Rbar' : ℝ) (E' : ℝ) (lower' : ℝ) (upper' : ℝ) (D' : ℝ),
    Galarza26RelativeCredit.selectionAndDomain qbar' a' lower' upper' ∧
      Galarza26RelativeCredit.technologyCapacityCost g' c' E' ∧
        Galarza26RelativeCredit.conditionedPayoff Rbar' lower' upper' ∧
          D' ∈ Set.Ioo lower' upper' ∧
            g' * Galarza26RelativeCredit.premium Rbar' D' / c' = E' ∧
              E' < Galarza26RelativeCredit.priorRisk qbar' a' D' / g' ∧
                ¬DifferentiableAt ℝ (Galarza26RelativeCredit.optimalRisk qbar' a' g' c' Rbar' E') D'
\end{ReviewVerbatim}

\paragraph{Retained governing declarations}
\begin{itemize}\raggedright
\item \hyperlink{governing-declaration-042b6300ffc11e3e}{Galarza26\allowbreak{}Relative\allowbreak{}Credit.\allowbreak{}Continuation\allowbreak{}Problem.\allowbreak{}Is\allowbreak{}Maximizer}
\item \hyperlink{governing-declaration-6e5d4449ad412574}{Galarza26\allowbreak{}Relative\allowbreak{}Credit.\allowbreak{}capacity\allowbreak{}Regime}
\item \hyperlink{paper-prerequisite-a6fae381e2d9a684}{Galarza26\allowbreak{}Relative\allowbreak{}Credit.\allowbreak{}conditioned\allowbreak{}Continuation\allowbreak{}Problem}
\item \hyperlink{paper-prerequisite-0bf7e0c122a51179}{Galarza26\allowbreak{}Relative\allowbreak{}Credit.\allowbreak{}conditioned\allowbreak{}Payoff}
\item \hyperlink{governing-declaration-65adc159d5a1934d}{Galarza26\allowbreak{}Relative\allowbreak{}Credit.\allowbreak{}floor\allowbreak{}Regime}
\item \hyperlink{governing-declaration-5f246d7c1bbfcd61}{Galarza26\allowbreak{}Relative\allowbreak{}Credit.\allowbreak{}interior\allowbreak{}Regime}
\item \hyperlink{governing-declaration-41e961e1975be029}{Galarza26\allowbreak{}Relative\allowbreak{}Credit.\allowbreak{}optimal\allowbreak{}Effort}
\item \hyperlink{governing-declaration-e139cf967ddac229}{Galarza26\allowbreak{}Relative\allowbreak{}Credit.\allowbreak{}optimal\allowbreak{}Risk}
\item \hyperlink{paper-prerequisite-6956aaccf5fffc8f}{Galarza26\allowbreak{}Relative\allowbreak{}Credit.\allowbreak{}premium}
\item \hyperlink{paper-prerequisite-e8657fe6a37f000a}{Galarza26\allowbreak{}Relative\allowbreak{}Credit.\allowbreak{}prior\allowbreak{}Risk}
\item \hyperlink{paper-prerequisite-e9553b2c455b69dd}{Galarza26\allowbreak{}Relative\allowbreak{}Credit.\allowbreak{}screened\allowbreak{}Risk}
\item \hyperlink{paper-prerequisite-ab0e2f9b84884bc6}{Galarza26\allowbreak{}Relative\allowbreak{}Credit.\allowbreak{}selection\allowbreak{}And\allowbreak{}Domain}
\item \hyperlink{paper-prerequisite-e5f9b7ce782f743a}{Galarza26\allowbreak{}Relative\allowbreak{}Credit.\allowbreak{}technology\allowbreak{}Capacity\allowbreak{}Cost}
\end{itemize}
\paragraph{Lean-elaborated claim roles}\begin{ReviewVerbatim}
1. conclusion: (∀ (qbar a g c Rbar F E lower upper : ℝ),
    Galarza26RelativeCredit.selectionAndDomain qbar a lower upper →
      Galarza26RelativeCredit.technologyCapacityCost g c E →
        Galarza26RelativeCredit.conditionedPayoff Rbar lower upper →
          (∀ D ∈ Set.Icc lower upper,
              ((∃! e : ℝ,
                    (Galarza26RelativeCredit.conditionedContinuationProblem qbar a g c Rbar F E D).IsMaximizer e) ∧
                  (Galarza26RelativeCredit.conditionedContinuationProblem qbar a g c Rbar F E D).IsMaximizer
                    (Galarza26RelativeCredit.optimalEffort qbar a g c Rbar E D)) ∧
                Galarza26RelativeCredit.screenedRisk qbar a g D
                      (Galarza26RelativeCredit.optimalEffort qbar a g c Rbar E D) =
                    Galarza26RelativeCredit.optimalRisk qbar a g c Rbar E D ∧
                  0 ≤ Galarza26RelativeCredit.optimalRisk qbar a g c Rbar E D ∧
                    Galarza26RelativeCredit.optimalRisk qbar a g c Rbar E D ≤ 1) ∧
            (∀ D ∈ Set.Ioo lower upper,
                Galarza26RelativeCredit.interiorRegime qbar a g c Rbar E D →
                  HasDerivAt (Galarza26RelativeCredit.optimalRisk qbar a g c Rbar E) (a - g ^ 2 / c) D) ∧
              (∀ D ∈ Set.Ioo lower upper,
                  Galarza26RelativeCredit.capacityRegime qbar a g c Rbar E D →
                    HasDerivAt (Galarza26RelativeCredit.optimalRisk qbar a g c Rbar E) a D) ∧
                (∀ D ∈ Set.Ioo lower upper,
                    Galarza26RelativeCredit.floorRegime qbar a g c Rbar E D →
                      HasDerivAt (Galarza26RelativeCredit.optimalRisk qbar a g c Rbar E) 0 D) ∧
                  (E = 0 →
                      (∀ D ∈ Set.Icc lower upper,
                          Galarza26RelativeCredit.optimalEffort qbar a g c Rbar E D = 0 ∧
                            Galarza26RelativeCredit.optimalRisk qbar a g c Rbar E D =
                              Galarza26RelativeCredit.priorRisk qbar a D) ∧
                        ∀ D ∈ Set.Ioo lower upper,
                          HasDerivAt (Galarza26RelativeCredit.optimalRisk qbar a g c Rbar E) a D) ∧
                    ContinuousOn (Galarza26RelativeCredit.optimalEffort qbar a g c Rbar E) (Set.Icc lower upper) ∧
                      ContinuousOn (Galarza26RelativeCredit.optimalRisk qbar a g c Rbar E) (Set.Icc lower upper) ∧
                        ∀ D ∈ Set.Icc lower upper,
                          ∃ (left : ℝ) (right : ℝ),
                            HasDerivWithinAt (Galarza26RelativeCredit.optimalRisk qbar a g c Rbar E) left (Set.Iic D)
                                D ∧
                              HasDerivWithinAt (Galarza26RelativeCredit.optimalRisk qbar a g c Rbar E) right (Set.Ici D)
                                D) ∧
  ∃ (qbar' : ℝ) (a' : ℝ) (g' : ℝ) (c' : ℝ) (Rbar' : ℝ) (E' : ℝ) (lower' : ℝ) (upper' : ℝ) (D' : ℝ),
    Galarza26RelativeCredit.selectionAndDomain qbar' a' lower' upper' ∧
      Galarza26RelativeCredit.technologyCapacityCost g' c' E' ∧
        Galarza26RelativeCredit.conditionedPayoff Rbar' lower' upper' ∧
          D' ∈ Set.Ioo lower' upper' ∧
            g' * Galarza26RelativeCredit.premium Rbar' D' / c' = E' ∧
              E' < Galarza26RelativeCredit.priorRisk qbar' a' D' / g' ∧
                ¬DifferentiableAt ℝ (Galarza26RelativeCredit.optimalRisk qbar' a' g' c' Rbar' E') D'
\end{ReviewVerbatim}

\paragraph{Lean proof endpoint (built separately)}\begin{ReviewVerbatim}
Galarza26RelativeCredit.proposition1
\end{ReviewVerbatim}

\paragraph{Recorded source-to-Spec screening}
\textbf{Verdict:} \texttt{matches}
\\
{\raggedright
\textbf{Reason:} Compared the verbatim Proposition 1 and Assumptions 1–3 with the expanded target and every displayed prerequisite.\allowbreak{} selectionAndDomain, technologyCapacityCost, and conditionedPayoff encode the stated compact probability domain, g > 0, c > 0, nonnegative capacity, and positive R(D);\allowbreak{} feasibleEfforts and bankProfit reproduce the feasible interval and objective (2).\allowbreak{} optimalEffort is min\{gR/\allowbreak{}c, E, q0/\allowbreak{}g\}, optimalRisk is max\{0, q0 -\allowbreak{} g min\{E, gR/\allowbreak{}c\}\}, and the three strict regime predicates yield exactly the slopes a -\allowbreak{} g\textasciicircum{}2/\allowbreak{}c, a, and 0.\allowbreak{} The target also retains the E = 0 case, continuity, and one-\allowbreak{}sided derivatives at constraint ties;\allowbreak{} its admissible tie witness for failure of an ordinary derivative is the source's stated possibility that such a derivative need not exist.\allowbreak{} The extra all-\allowbreak{}domain bounds and one-\allowbreak{}sided statements are valid strengthenings that include every source-\allowbreak{}admissible case and do not alter the source conclusion.\allowbreak{}
\par}
\reviewmatch{Matches source input}
\noindent\textbf{Reviewer annotation}\par
\reviewerbox
\clearpage
\hypertarget{source-claim-f0d8e27dfc3b3ee3}{}
\section*{Review row 2}
\textbf{Source locator:} {\footnotesize\raggedright cited publication, lines 484-\allowbreak{}485\par}
\paragraph{Verbatim source input}
\begin{ReviewVerbatim}
Proposición 2 (Umbral condicional). Bajo los Supuestos 1–4, existe un único κ∗ ∈ (0, 1) con
∂D q ∗ (D0 , κ∗ ) = 0. La derivada es positiva para κ < κ∗ y negativa para κ > κ∗ .

[Next verbatim source excerpt]

primarias consultadas. No es evidencia de que nuestro indicador funcione. Su construcción separa
el deterioro de créditos inicialmente normales durante una ventana de doce meses del ratio contable
de morosidad. Esta diferencia respalda la necesidad de explicitar evento y horizonte; no justifica
equiparar automáticamente nuestro q ∗ con una serie publicada. Clark and West (2007) proporciona
una comparación de modelos predictivos anidados, relevante para un trabajo empírico posterior.
No se calcula aquí una prueba ni un valor p.
Los mecanismos rivales tienen implicaciones que ayudan a diseñar esa prueba. Si domina selección, precios altos anteceden peor calidad de nuevas cosechas, manteniendo constante la evaluación.
Si domina evaluación, el precio alto debe acompañarse de mayor esfuerzo efectivo o mejor documentación y menor riesgo posterior, salvo que la capacidad esté saturada. Si domina composición, el
cambio de signo aparece al agregar al saldo y depende de la participación de colocaciones recientes.
Los costos de fondeo o venta cruzada pueden generar brechas sin estos patrones. Con promedios
públicos será difícil separarlos; harían falta, idealmente, datos de cosechas, rotación, evaluación y
garantías.

4.

Ambiente y decisión del banco

4.1.

Agentes, información, unidades y cronología

El único agente que optimiza en este modelo es un banco neutral al riesgo. Los solicitantes

[Next verbatim source excerpt]

y los bancos pares determinan una calidad y una referencia reducidas; no se modela su juego
estratégico. El supervisor observa D y más tarde un ratio de cartera, pero no el esfuerzo ni toda la

[Next verbatim source excerpt]

conjunto de solicitantes y el banco observa su riesgo medio; luego elige esfuerzo antes de originar;
finalmente se realiza el estado crediticio y se observa la cartera al horizonte h. El banco conoce los
parámetros al decidir. La incertidumbre de desempeño permanece y se integra en su objetivo. El
precio es un argumento del problema de continuación, no una variable resuelta en equilibrio.

[Next verbatim source excerpt]

problemas con precios cercanos, dejando fijo el segmento y permitiendo que cambie el esfuerzo
óptimo. Solo en la Proposición 2 se comparan distintos valores de κ, siempre al mismo precio D0 .
La prima de M ′ se refiere a D; las primas de g ′ y c′ en el análisis de opacidad se refieren a κ.
La tabla siguiente sirve como diccionario. Los objetos de cartera se usan recién en la Proposición 3; no intervienen en la elección inicial del esfuerzo.
Símbolo

Significado

Unidad o restricción

D
κ
q̄, a
q 0 , q, q ∗

[Next verbatim source excerpt]

Esfuerzo, capacidad y máximo esfuerzo factible
Productividad y escala del costo de evaluación
Premio por desempeño, premio de referencia y costo fijo
Esfuerzo deseado sin restricciones: gR/c
Fuerza de evaluación g 2 /c y pendiente neta a − H
Derivada de q ∗ en el régimen efectivo
Saldos antiguo, nuevo y nuevo cuando D = 0
Atraso antiguo, peso nuevo y ratio total
Sensibilidad proporcional de B al precio

Tasa anual decimal
Índice en [0, 1]
Probabilidad; a > 0
Mismo evento y horizonte
Índices no negativos
Ambos positivos
Pagos por sol de crédito
Índice de esfuerzo
Riesgo por unidad de D
Riesgo por unidad de D
Misma unidad monetaria
Fracciones en [0, 1]
Por unidad de D

Supuesto 1 (Selección y dominio). Para cada κ fijo, a = a(κ) > 0, q̄ ∈ (0, 1) y D pertenece a un
intervalo compacto D. La probabilidad previa es q 0 (D, κ) = q̄ + a(κ)D ∈ [0, 1] en todo D.
La forma afín es una aproximación local al cambio de solicitantes. Una subida de un punto
porcentual es ∆D = 0.01: cambia la probabilidad en 0.01a unidades, equivalentes a a puntos porcentuales. Por ejemplo, con a = 0.18 el cambio es 0.0018 de probabilidad, es decir, 0.18 puntos
porcentuales. No se extrapola fuera del dominio admisible. La pendiente positiva representa selección adversa, no una ley empírica universal ni una respuesta derivada de un problema explícito de
cada deudor.
La restricción de probabilidad también fija qué precios se pueden comparar. Restar q̄ y dividir
por a > 0 da
1 − q̄
q̄
.
0 ≤ q̄ + aD ≤ 1 ⇐⇒ − ≤ D ≤
a
a
El dominio compacto D se elige dentro de ese intervalo y debe respetar además el pago positivo del
Supuesto 3. No se pretende describir cualquier tasa posible con una probabilidad afín.

[Next verbatim source excerpt]

no está resuelta: el resultado se refiere a una unidad que el banco ya contempla originar. La remuneración supone que evitar el evento adverso mejora el pago esperado. Recuperaciones, seguros y

[Next verbatim source excerpt]

garantías que varían con precio o esfuerzo exigirían modificar este premio. El horizonte del evento
y el flujo contractual deben ser compatibles en una implementación empírica.
Para leer el objetivo, empiece por 1 − q 0 + ge = 1 − q: es la probabilidad de evitar el evento
adverso. Multiplicarla por R da el premio esperado. Después se restan el costo fijo F y el costo de

[Next verbatim source excerpt]

de morosidad. Esta diferencia respalda la necesidad de explicitar evento y horizonte; no justifica
equiparar automáticamente nuestro q ∗ con una serie publicada. Clark and West (2007) proporciona
una comparación de modelos predictivos anidados, relevante para un trabajo empírico posterior.
No se calcula aquí una prueba ni un valor p.
Los mecanismos rivales tienen implicaciones que ayudan a diseñar esa prueba. Si domina selección, precios altos anteceden peor calidad de nuevas cosechas, manteniendo constante la evaluación.
Si domina evaluación, el precio alto debe acompañarse de mayor esfuerzo efectivo o mejor documentación y menor riesgo posterior, salvo que la capacidad esté saturada. Si domina composición, el
cambio de signo aparece al agregar al saldo y depende de la participación de colocaciones recientes.
Los costos de fondeo o venta cruzada pueden generar brechas sin estos patrones. Con promedios
públicos será difícil separarlos; harían falta, idealmente, datos de cosechas, rotación, evaluación y
garantías.

4.

Ambiente y decisión del banco

4.1.

Agentes, información, unidades y cronología

El único agente que optimiza en este modelo es un banco neutral al riesgo. Los solicitantes
y los bancos pares determinan una calidad y una referencia reducidas; no se modela su juego
estratégico. El supervisor observa D y más tarde un ratio de cartera, pero no el esfuerzo ni toda la

[Next verbatim source excerpt]

composición de riesgo. La opacidad κ ∈ [0, 1] es una característica exógena del segmento: un valor
mayor representa menor información verificable inicial. No se estima este índice ni se equipara

[Next verbatim source excerpt]

parámetros al decidir. La incertidumbre de desempeño permanece y se integra en su objetivo. El
precio es un argumento del problema de continuación, no una variable resuelta en equilibrio.
q 0 y q son probabilidades sin unidades, referidas al mismo evento y horizonte para una cohorte
estandarizada. e y κ son índices. D se mide en unidades de tasa anual decimal. Los pagos y costos

[Next verbatim source excerpt]

La prima de M ′ se refiere a D; las primas de g ′ y c′ en el análisis de opacidad se refieren a κ.
La tabla siguiente sirve como diccionario. Los objetos de cartera se usan recién en la Proposición 3; no intervienen en la elección inicial del esfuerzo.
Símbolo

Significado

Unidad o restricción

D
κ
q̄, a
q 0 , q, q ∗
e, ē, u
g, c

[Next verbatim source excerpt]

Premio por desempeño, premio de referencia y costo fijo
Esfuerzo deseado sin restricciones: gR/c
Fuerza de evaluación g 2 /c y pendiente neta a − H
Derivada de q ∗ en el régimen efectivo
Saldos antiguo, nuevo y nuevo cuando D = 0
Atraso antiguo, peso nuevo y ratio total
Sensibilidad proporcional de B al precio

Tasa anual decimal
Índice en [0, 1]
Probabilidad; a > 0
Mismo evento y horizonte
Índices no negativos
Ambos positivos
Pagos por sol de crédito
Índice de esfuerzo
Riesgo por unidad de D
Riesgo por unidad de D
Misma unidad monetaria
Fracciones en [0, 1]
Por unidad de D

Supuesto 1 (Selección y dominio). Para cada κ fijo, a = a(κ) > 0, q̄ ∈ (0, 1) y D pertenece a un
intervalo compacto D. La probabilidad previa es q 0 (D, κ) = q̄ + a(κ)D ∈ [0, 1] en todo D.
La forma afín es una aproximación local al cambio de solicitantes. Una subida de un punto
porcentual es ∆D = 0.01: cambia la probabilidad en 0.01a unidades, equivalentes a a puntos porcentuales. Por ejemplo, con a = 0.18 el cambio es 0.0018 de probabilidad, es decir, 0.18 puntos
porcentuales. No se extrapola fuera del dominio admisible. La pendiente positiva representa selección adversa, no una ley empírica universal ni una respuesta derivada de un problema explícito de
cada deudor.
La restricción de probabilidad también fija qué precios se pueden comparar. Restar q̄ y dividir
por a > 0 da
1 − q̄
q̄
.
0 ≤ q̄ + aD ≤ 1 ⇐⇒ − ≤ D ≤
a
a
El dominio compacto D se elige dentro de ese intervalo y debe respetar además el pago positivo del

[Next verbatim source excerpt]

todo el objetivo tampoco lo afecta si no depende de e. La participación o decisión de no prestar

[Next verbatim source excerpt]

no está resuelta: el resultado se refiere a una unidad que el banco ya contempla originar. La remuneración supone que evitar el evento adverso mejora el pago esperado. Recuperaciones, seguros y
garantías que varían con precio o esfuerzo exigirían modificar este premio. El horizonte del evento
y el flujo contractual deben ser compatibles en una implementación empírica.
Para leer el objetivo, empiece por 1 − q 0 + ge = 1 − q: es la probabilidad de evitar el evento

[Next verbatim source excerpt]

El único agente que optimiza en este modelo es un banco neutral al riesgo. Los solicitantes
y los bancos pares determinan una calidad y una referencia reducidas; no se modela su juego
estratégico. El supervisor observa D y más tarde un ratio de cartera, pero no el esfuerzo ni toda la
composición de riesgo. La opacidad κ ∈ [0, 1] es una característica exógena del segmento: un valor
mayor representa menor información verificable inicial. No se estima este índice ni se equipara
automáticamente microempresa con opacidad alta.
La cronología es: primero se fijan segmento, tecnología y precio observado; después se forma el
conjunto de solicitantes y el banco observa su riesgo medio; luego elige esfuerzo antes de originar;
finalmente se realiza el estado crediticio y se observa la cartera al horizonte h. El banco conoce los
parámetros al decidir. La incertidumbre de desempeño permanece y se integra en su objetivo. El
precio es un argumento del problema de continuación, no una variable resuelta en equilibrio.
q 0 y q son probabilidades sin unidades, referidas al mismo evento y horizonte para una cohorte
estandarizada. e y κ son índices. D se mide en unidades de tasa anual decimal. Los pagos y costos
se normalizan por sol de crédito. La sensibilidad monetaria del premio de desempeño a una unidad
4


[form-feed in source extraction]
de D se normaliza a uno. Esto es una simplificación de remuneración, no una conversión exacta de
tasas efectivas a cualquier plazo contractual.
Conviene separar tres operaciones que aparecerán en las pruebas. Al elegir e, el banco mantiene
fijos D y κ: compara esfuerzos para un mismo precio y segmento. Al calcular ∂D q ∗ , compara
problemas con precios cercanos, dejando fijo el segmento y permitiendo que cambie el esfuerzo
óptimo. Solo en la Proposición 2 se comparan distintos valores de κ, siempre al mismo precio D0 .
La prima de M ′ se refiere a D; las primas de g ′ y c′ en el análisis de opacidad se refieren a κ.

[Next verbatim source excerpt]

Supuesto 1 (Selección y dominio). Para cada κ fijo, a = a(κ) > 0, q̄ ∈ (0, 1) y D pertenece a un
intervalo compacto D. La probabilidad previa es q 0 (D, κ) = q̄ + a(κ)D ∈ [0, 1] en todo D.

[Next verbatim source excerpt]

Supuesto 2 (Tecnología, capacidad y costo). g = g(κ) > 0, c = c(κ) > 0 y ē ≥ 0 son finitos.
Para un volumen normalizado de créditos, q(D, e, κ) = q 0 (D, κ) − g(κ)e y el costo es c(κ)e2 /2. La
elección factible cumple
0 ≤ e ≤ u(D, κ) := mı́n{ē, q 0 (D, κ)/g(κ)}.

[Next verbatim source excerpt]

Supuesto 3 (Remuneración y problema condicionado). El premio normalizado por desempeño es
R(D) = R̄ + D > 0 y F es un costo por unidad independiente de e. El banco resuelve
máx
e∈[0,u(D,κ)]

π(e; D, κ) = [1 − q 0 (D, κ) + g(κ)e]R(D) − F −

c(κ) 2
e .
2

(2)

El precio, la referencia de pares y κ se mantienen fijos al evaluar desviaciones de esfuerzo.
El costo F no afecta el óptimo de esfuerzo. Un volumen estrictamente positivo que multiplique
todo el objetivo tampoco lo afecta si no depende de e. La participación o decisión de no prestar
no está resuelta: el resultado se refiere a una unidad que el banco ya contempla originar. La remuneración supone que evitar el evento adverso mejora el pago esperado. Recuperaciones, seguros y
garantías que varían con precio o esfuerzo exigirían modificar este premio. El horizonte del evento
y el flujo contractual deben ser compatibles en una implementación empírica.

[Next verbatim source excerpt]

Proposición 1 (Solución global y pendientes por régimen). Bajo los Supuestos 1–3, existe un
único máximo en (2). Es
[U+001A control character]
[U+001B control character]
[U+001A control character]
[U+001A control character]
[U+001B control character][U+001B control character]
gR(D)
q 0 (D)
gR(D)
∗
0
∗
e = mı́n
, ē,
,
q = máx 0, q (D) − g mı́n ē,
.
(3)
c
g
c
Para κ fijo, en cualquier intervalo abierto con desigualdades estrictas que mantengan el régimen,
las pendientes son
[U+F8F1 delimiter glyph]
[U+F8F4 delimiter glyph]a − g 2 /c, gR/c < ē y gR/c < q 0 /g,
∂q ∗ [U+F8F2 delimiter glyph]
= a,
(4)
ē < gR/c y ē < q 0 /g,
[U+F8F4 delimiter glyph]
∂D
[U+F8F3 delimiter glyph]
0
0
0,
q /g < gR/c y q /g < ē.
Si ē = 0, e∗ = 0 y q ∗ = q 0 , con pendiente a en el interior del dominio de D. En empates de
restricciones la solución es continua y se usan derivadas laterales; una derivada ordinaria puede
no existir.

[Next verbatim source excerpt]

Defina H(κ) = g(κ)2 /c(κ) y β(κ) = a(κ) − H(κ).

[Next verbatim source excerpt]

Supuesto 4 (Condiciones suficientes para el cruce). En κ ∈ [0, 1], a, g, c son continuas y positivas.
La función a − H es estrictamente decreciente y satisface a(0) > H(0) y a(1) < H(1). Al precio D0
que se compara se cumplen estrictamente las dos condiciones de interioridad de la primera línea de
(4) para todo κ.
\end{ReviewVerbatim}



\paragraph{Expanded PaperInterface specification}
\begin{ReviewVerbatim}
∀ (qbar Rbar F E lower upper D0 : ℝ) (a g c : ℝ → ℝ),
  D0 ∈ Set.Icc lower upper →
    (∀ k ∈ Set.Icc 0 1,
        Galarza26RelativeCredit.selectionAndDomain qbar (a k) lower upper ∧
          Galarza26RelativeCredit.technologyCapacityCost (g k) (c k) E ∧
            Galarza26RelativeCredit.conditionedPayoff Rbar lower upper) →
      Galarza26RelativeCredit.crossingConditions qbar Rbar E D0 a g c →
        ∃ k ∈ Set.Ioo 0 1,
          HasDerivAt (Galarza26RelativeCredit.optimalRisk qbar (a k) (g k) (c k) Rbar E) 0 D0 ∧
            (∀ y ∈ Set.Ioo 0 1,
                HasDerivAt (Galarza26RelativeCredit.optimalRisk qbar (a y) (g y) (c y) Rbar E) 0 D0 → y = k) ∧
              (∀ x ∈ Set.Icc 0 1,
                  x < k →
                    0 < Galarza26RelativeCredit.screeningAdvantage (a x) (g x) (c x) ∧
                      HasDerivAt (Galarza26RelativeCredit.optimalRisk qbar (a x) (g x) (c x) Rbar E)
                        (Galarza26RelativeCredit.screeningAdvantage (a x) (g x) (c x)) D0) ∧
                ∀ x ∈ Set.Icc 0 1,
                  k < x →
                    Galarza26RelativeCredit.screeningAdvantage (a x) (g x) (c x) < 0 ∧
                      HasDerivAt (Galarza26RelativeCredit.optimalRisk qbar (a x) (g x) (c x) Rbar E)
                        (Galarza26RelativeCredit.screeningAdvantage (a x) (g x) (c x)) D0
\end{ReviewVerbatim}

\paragraph{Retained governing declarations}
\begin{itemize}\raggedright
\item \hyperlink{paper-prerequisite-0bf7e0c122a51179}{Galarza26\allowbreak{}Relative\allowbreak{}Credit.\allowbreak{}conditioned\allowbreak{}Payoff}
\item \hyperlink{paper-prerequisite-4f59b48de0be00ac}{Galarza26\allowbreak{}Relative\allowbreak{}Credit.\allowbreak{}crossing\allowbreak{}Conditions}
\item \hyperlink{governing-declaration-e139cf967ddac229}{Galarza26\allowbreak{}Relative\allowbreak{}Credit.\allowbreak{}optimal\allowbreak{}Risk}
\item \hyperlink{paper-prerequisite-29fcb63a4a0578b3}{Galarza26\allowbreak{}Relative\allowbreak{}Credit.\allowbreak{}screening\allowbreak{}Advantage}
\item \hyperlink{paper-prerequisite-ab0e2f9b84884bc6}{Galarza26\allowbreak{}Relative\allowbreak{}Credit.\allowbreak{}selection\allowbreak{}And\allowbreak{}Domain}
\item \hyperlink{paper-prerequisite-e5f9b7ce782f743a}{Galarza26\allowbreak{}Relative\allowbreak{}Credit.\allowbreak{}technology\allowbreak{}Capacity\allowbreak{}Cost}
\end{itemize}
\paragraph{Lean-elaborated claim roles}\begin{ReviewVerbatim}
1. parameter: ℝ
2. parameter: ℝ
3. parameter: ℝ
4. parameter: ℝ
5. parameter: ℝ
6. parameter: ℝ
7. parameter: ℝ
8. parameter: ℝ → ℝ
9. parameter: ℝ → ℝ
10. parameter: ℝ → ℝ
11. assumption: D0 ∈ Set.Icc lower upper
12. assumption: ∀ k ∈ Set.Icc 0 1,
  Galarza26RelativeCredit.selectionAndDomain qbar (a k) lower upper ∧
    Galarza26RelativeCredit.technologyCapacityCost (g k) (c k) E ∧
      Galarza26RelativeCredit.conditionedPayoff Rbar lower upper
13. assumption: Galarza26RelativeCredit.crossingConditions qbar Rbar E D0 a g c
14. conclusion: ∃ k ∈ Set.Ioo 0 1,
  HasDerivAt (Galarza26RelativeCredit.optimalRisk qbar (a k) (g k) (c k) Rbar E) 0 D0 ∧
    (∀ y ∈ Set.Ioo 0 1, HasDerivAt (Galarza26RelativeCredit.optimalRisk qbar (a y) (g y) (c y) Rbar E) 0 D0 → y = k) ∧
      (∀ x ∈ Set.Icc 0 1,
          x < k →
            0 < Galarza26RelativeCredit.screeningAdvantage (a x) (g x) (c x) ∧
              HasDerivAt (Galarza26RelativeCredit.optimalRisk qbar (a x) (g x) (c x) Rbar E)
                (Galarza26RelativeCredit.screeningAdvantage (a x) (g x) (c x)) D0) ∧
        ∀ x ∈ Set.Icc 0 1,
          k < x →
            Galarza26RelativeCredit.screeningAdvantage (a x) (g x) (c x) < 0 ∧
              HasDerivAt (Galarza26RelativeCredit.optimalRisk qbar (a x) (g x) (c x) Rbar E)
                (Galarza26RelativeCredit.screeningAdvantage (a x) (g x) (c x)) D0
\end{ReviewVerbatim}

\paragraph{Lean proof endpoint (built separately)}\begin{ReviewVerbatim}
Galarza26RelativeCredit.proposition2
\end{ReviewVerbatim}

\paragraph{Recorded source-to-Spec screening}
\textbf{Verdict:} \texttt{matches}
\\
{\raggedright
\textbf{Reason:} Compared the verbatim Proposition 2 and Assumptions 1–4 with the full expanded target and prerequisites.\allowbreak{} For every kappa in [0,1], the displayed assumptions impose the source's probability domain, positive technology parameters, nonnegative capacity, and positive payoff.\allowbreak{} crossingConditions states continuity and positivity of a, g, c, strict decrease of beta(kappa)=a(kappa)-\allowbreak{}g(kappa)\textasciicircum{}2/\allowbreak{}c(kappa), the positive and negative endpoint signs, and both strict interior-\allowbreak{}regime inequalities at the common D0.\allowbreak{} The conclusion then gives a unique interior kappa with D-\allowbreak{}derivative zero and, on the two sides, the exact positive/\allowbreak{}negative beta derivative required by the source.\allowbreak{} Quantifiers, the fixed-\allowbreak{}price comparison, strict inequalities, and the interval domains agree;\allowbreak{} the unused fixed cost F does not change the source claim.\allowbreak{}
\par}
\reviewmatch{Matches source input}
\noindent\textbf{Reviewer annotation}\par
\reviewerbox
\clearpage
\hypertarget{source-claim-194f750a63a93998}{}
\section*{Review row 3}
\textbf{Source locator:} {\footnotesize\raggedright cited publication, lines 540-\allowbreak{}551\par}
\paragraph{Verbatim source input}
\begin{ReviewVerbatim}
Proposición 3 (Transmisión y composición). Bajo los Supuestos 1–3 y 5, el ratio esperado de
cartera es M (D) = (1 − w)m0 + wq ∗ (D). En cualquier régimen diferenciable, escribiendo βD =
∂D q ∗ (D),
M ′ (D) = wβD + w′ (q ∗ − m0 ).
(6)
Si B es constante, el signo es el de βD . Si βD ̸= 0 y |w′ (q ∗ − m0 )| < w|βD |, el signo también se
conserva. Con B(D) = B1 e−εD , B1 > 0, ε ≥ 0,
M ′ = w{βD − ε(1 − w)(q ∗ − m0 )}.

(7)

En particular, βD > 0 y q ∗ > m0 generan M ′ < 0 si y solo si ε(1 − w)(q ∗ − m0 ) > βD .

[Next verbatim source excerpt]

primarias consultadas. No es evidencia de que nuestro indicador funcione. Su construcción separa
el deterioro de créditos inicialmente normales durante una ventana de doce meses del ratio contable
de morosidad. Esta diferencia respalda la necesidad de explicitar evento y horizonte; no justifica
equiparar automáticamente nuestro q ∗ con una serie publicada. Clark and West (2007) proporciona
una comparación de modelos predictivos anidados, relevante para un trabajo empírico posterior.
No se calcula aquí una prueba ni un valor p.
Los mecanismos rivales tienen implicaciones que ayudan a diseñar esa prueba. Si domina selección, precios altos anteceden peor calidad de nuevas cosechas, manteniendo constante la evaluación.
Si domina evaluación, el precio alto debe acompañarse de mayor esfuerzo efectivo o mejor documentación y menor riesgo posterior, salvo que la capacidad esté saturada. Si domina composición, el
cambio de signo aparece al agregar al saldo y depende de la participación de colocaciones recientes.
Los costos de fondeo o venta cruzada pueden generar brechas sin estos patrones. Con promedios
públicos será difícil separarlos; harían falta, idealmente, datos de cosechas, rotación, evaluación y
garantías.

4.

Ambiente y decisión del banco

4.1.

Agentes, información, unidades y cronología

El único agente que optimiza en este modelo es un banco neutral al riesgo. Los solicitantes

[Next verbatim source excerpt]

y los bancos pares determinan una calidad y una referencia reducidas; no se modela su juego
estratégico. El supervisor observa D y más tarde un ratio de cartera, pero no el esfuerzo ni toda la

[Next verbatim source excerpt]

conjunto de solicitantes y el banco observa su riesgo medio; luego elige esfuerzo antes de originar;
finalmente se realiza el estado crediticio y se observa la cartera al horizonte h. El banco conoce los
parámetros al decidir. La incertidumbre de desempeño permanece y se integra en su objetivo. El
precio es un argumento del problema de continuación, no una variable resuelta en equilibrio.

[Next verbatim source excerpt]

problemas con precios cercanos, dejando fijo el segmento y permitiendo que cambie el esfuerzo
óptimo. Solo en la Proposición 2 se comparan distintos valores de κ, siempre al mismo precio D0 .
La prima de M ′ se refiere a D; las primas de g ′ y c′ en el análisis de opacidad se refieren a κ.
La tabla siguiente sirve como diccionario. Los objetos de cartera se usan recién en la Proposición 3; no intervienen en la elección inicial del esfuerzo.
Símbolo

Significado

Unidad o restricción

D
κ
q̄, a
q 0 , q, q ∗

[Next verbatim source excerpt]

Esfuerzo, capacidad y máximo esfuerzo factible
Productividad y escala del costo de evaluación
Premio por desempeño, premio de referencia y costo fijo
Esfuerzo deseado sin restricciones: gR/c
Fuerza de evaluación g 2 /c y pendiente neta a − H
Derivada de q ∗ en el régimen efectivo
Saldos antiguo, nuevo y nuevo cuando D = 0
Atraso antiguo, peso nuevo y ratio total
Sensibilidad proporcional de B al precio

Tasa anual decimal
Índice en [0, 1]
Probabilidad; a > 0
Mismo evento y horizonte
Índices no negativos
Ambos positivos
Pagos por sol de crédito
Índice de esfuerzo
Riesgo por unidad de D
Riesgo por unidad de D
Misma unidad monetaria
Fracciones en [0, 1]
Por unidad de D

Supuesto 1 (Selección y dominio). Para cada κ fijo, a = a(κ) > 0, q̄ ∈ (0, 1) y D pertenece a un
intervalo compacto D. La probabilidad previa es q 0 (D, κ) = q̄ + a(κ)D ∈ [0, 1] en todo D.
La forma afín es una aproximación local al cambio de solicitantes. Una subida de un punto
porcentual es ∆D = 0.01: cambia la probabilidad en 0.01a unidades, equivalentes a a puntos porcentuales. Por ejemplo, con a = 0.18 el cambio es 0.0018 de probabilidad, es decir, 0.18 puntos
porcentuales. No se extrapola fuera del dominio admisible. La pendiente positiva representa selección adversa, no una ley empírica universal ni una respuesta derivada de un problema explícito de
cada deudor.
La restricción de probabilidad también fija qué precios se pueden comparar. Restar q̄ y dividir
por a > 0 da
1 − q̄
q̄
.
0 ≤ q̄ + aD ≤ 1 ⇐⇒ − ≤ D ≤
a
a
El dominio compacto D se elige dentro de ese intervalo y debe respetar además el pago positivo del
Supuesto 3. No se pretende describir cualquier tasa posible con una probabilidad afín.

[Next verbatim source excerpt]

q = máx 0, q (D) − g mı́n ē,
.
(3)

[Next verbatim source excerpt]

Para κ fijo, en cualquier intervalo abierto con desigualdades estrictas que mantengan el régimen,
las pendientes son
[U+F8F1 delimiter glyph]
[U+F8F4 delimiter glyph]a − g 2 /c, gR/c < ē y gR/c < q 0 /g,
∂q ∗ [U+F8F2 delimiter glyph]
= a,

[Next verbatim source excerpt]

de morosidad. Esta diferencia respalda la necesidad de explicitar evento y horizonte; no justifica
equiparar automáticamente nuestro q ∗ con una serie publicada. Clark and West (2007) proporciona
una comparación de modelos predictivos anidados, relevante para un trabajo empírico posterior.
No se calcula aquí una prueba ni un valor p.
Los mecanismos rivales tienen implicaciones que ayudan a diseñar esa prueba. Si domina selección, precios altos anteceden peor calidad de nuevas cosechas, manteniendo constante la evaluación.
Si domina evaluación, el precio alto debe acompañarse de mayor esfuerzo efectivo o mejor documentación y menor riesgo posterior, salvo que la capacidad esté saturada. Si domina composición, el
cambio de signo aparece al agregar al saldo y depende de la participación de colocaciones recientes.
Los costos de fondeo o venta cruzada pueden generar brechas sin estos patrones. Con promedios
públicos será difícil separarlos; harían falta, idealmente, datos de cosechas, rotación, evaluación y
garantías.

4.

Ambiente y decisión del banco

4.1.

Agentes, información, unidades y cronología

El único agente que optimiza en este modelo es un banco neutral al riesgo. Los solicitantes
y los bancos pares determinan una calidad y una referencia reducidas; no se modela su juego
estratégico. El supervisor observa D y más tarde un ratio de cartera, pero no el esfuerzo ni toda la

[Next verbatim source excerpt]

composición de riesgo. La opacidad κ ∈ [0, 1] es una característica exógena del segmento: un valor
mayor representa menor información verificable inicial. No se estima este índice ni se equipara

[Next verbatim source excerpt]

parámetros al decidir. La incertidumbre de desempeño permanece y se integra en su objetivo. El
precio es un argumento del problema de continuación, no una variable resuelta en equilibrio.
q 0 y q son probabilidades sin unidades, referidas al mismo evento y horizonte para una cohorte
estandarizada. e y κ son índices. D se mide en unidades de tasa anual decimal. Los pagos y costos

[Next verbatim source excerpt]

La prima de M ′ se refiere a D; las primas de g ′ y c′ en el análisis de opacidad se refieren a κ.
La tabla siguiente sirve como diccionario. Los objetos de cartera se usan recién en la Proposición 3; no intervienen en la elección inicial del esfuerzo.
Símbolo

Significado

Unidad o restricción

D
κ
q̄, a
q 0 , q, q ∗
e, ē, u
g, c

[Next verbatim source excerpt]

Premio por desempeño, premio de referencia y costo fijo
Esfuerzo deseado sin restricciones: gR/c
Fuerza de evaluación g 2 /c y pendiente neta a − H
Derivada de q ∗ en el régimen efectivo
Saldos antiguo, nuevo y nuevo cuando D = 0
Atraso antiguo, peso nuevo y ratio total
Sensibilidad proporcional de B al precio

Tasa anual decimal
Índice en [0, 1]
Probabilidad; a > 0
Mismo evento y horizonte
Índices no negativos
Ambos positivos
Pagos por sol de crédito
Índice de esfuerzo
Riesgo por unidad de D
Riesgo por unidad de D
Misma unidad monetaria
Fracciones en [0, 1]
Por unidad de D

Supuesto 1 (Selección y dominio). Para cada κ fijo, a = a(κ) > 0, q̄ ∈ (0, 1) y D pertenece a un
intervalo compacto D. La probabilidad previa es q 0 (D, κ) = q̄ + a(κ)D ∈ [0, 1] en todo D.
La forma afín es una aproximación local al cambio de solicitantes. Una subida de un punto
porcentual es ∆D = 0.01: cambia la probabilidad en 0.01a unidades, equivalentes a a puntos porcentuales. Por ejemplo, con a = 0.18 el cambio es 0.0018 de probabilidad, es decir, 0.18 puntos
porcentuales. No se extrapola fuera del dominio admisible. La pendiente positiva representa selección adversa, no una ley empírica universal ni una respuesta derivada de un problema explícito de
cada deudor.
La restricción de probabilidad también fija qué precios se pueden comparar. Restar q̄ y dividir
por a > 0 da
1 − q̄
q̄
.
0 ≤ q̄ + aD ≤ 1 ⇐⇒ − ≤ D ≤
a
a
El dominio compacto D se elige dentro de ese intervalo y debe respetar además el pago positivo del

[Next verbatim source excerpt]

,
q = máx 0, q (D) − g mı́n ē,
.

[Next verbatim source excerpt]

g
c
Para κ fijo, en cualquier intervalo abierto con desigualdades estrictas que mantengan el régimen,
las pendientes son
[U+F8F1 delimiter glyph]
[U+F8F4 delimiter glyph]a − g 2 /c, gR/c < ē y gR/c < q 0 /g,

[Next verbatim source excerpt]

El único agente que optimiza en este modelo es un banco neutral al riesgo. Los solicitantes
y los bancos pares determinan una calidad y una referencia reducidas; no se modela su juego
estratégico. El supervisor observa D y más tarde un ratio de cartera, pero no el esfuerzo ni toda la
composición de riesgo. La opacidad κ ∈ [0, 1] es una característica exógena del segmento: un valor
mayor representa menor información verificable inicial. No se estima este índice ni se equipara
automáticamente microempresa con opacidad alta.
La cronología es: primero se fijan segmento, tecnología y precio observado; después se forma el
conjunto de solicitantes y el banco observa su riesgo medio; luego elige esfuerzo antes de originar;
finalmente se realiza el estado crediticio y se observa la cartera al horizonte h. El banco conoce los
parámetros al decidir. La incertidumbre de desempeño permanece y se integra en su objetivo. El
precio es un argumento del problema de continuación, no una variable resuelta en equilibrio.
q 0 y q son probabilidades sin unidades, referidas al mismo evento y horizonte para una cohorte
estandarizada. e y κ son índices. D se mide en unidades de tasa anual decimal. Los pagos y costos
se normalizan por sol de crédito. La sensibilidad monetaria del premio de desempeño a una unidad
4


[form-feed in source extraction]
de D se normaliza a uno. Esto es una simplificación de remuneración, no una conversión exacta de
tasas efectivas a cualquier plazo contractual.
Conviene separar tres operaciones que aparecerán en las pruebas. Al elegir e, el banco mantiene
fijos D y κ: compara esfuerzos para un mismo precio y segmento. Al calcular ∂D q ∗ , compara
problemas con precios cercanos, dejando fijo el segmento y permitiendo que cambie el esfuerzo
óptimo. Solo en la Proposición 2 se comparan distintos valores de κ, siempre al mismo precio D0 .
La prima de M ′ se refiere a D; las primas de g ′ y c′ en el análisis de opacidad se refieren a κ.

[Next verbatim source excerpt]

Supuesto 1 (Selección y dominio). Para cada κ fijo, a = a(κ) > 0, q̄ ∈ (0, 1) y D pertenece a un
intervalo compacto D. La probabilidad previa es q 0 (D, κ) = q̄ + a(κ)D ∈ [0, 1] en todo D.

[Next verbatim source excerpt]

Supuesto 2 (Tecnología, capacidad y costo). g = g(κ) > 0, c = c(κ) > 0 y ē ≥ 0 son finitos.
Para un volumen normalizado de créditos, q(D, e, κ) = q 0 (D, κ) − g(κ)e y el costo es c(κ)e2 /2. La
elección factible cumple
0 ≤ e ≤ u(D, κ) := mı́n{ē, q 0 (D, κ)/g(κ)}.

[Next verbatim source excerpt]

Supuesto 3 (Remuneración y problema condicionado). El premio normalizado por desempeño es
R(D) = R̄ + D > 0 y F es un costo por unidad independiente de e. El banco resuelve
máx
e∈[0,u(D,κ)]

π(e; D, κ) = [1 − q 0 (D, κ) + g(κ)e]R(D) − F −

c(κ) 2
e .
2

(2)

El precio, la referencia de pares y κ se mantienen fijos al evaluar desviaciones de esfuerzo.
El costo F no afecta el óptimo de esfuerzo. Un volumen estrictamente positivo que multiplique
todo el objetivo tampoco lo afecta si no depende de e. La participación o decisión de no prestar
no está resuelta: el resultado se refiere a una unidad que el banco ya contempla originar. La remuneración supone que evitar el evento adverso mejora el pago esperado. Recuperaciones, seguros y
garantías que varían con precio o esfuerzo exigirían modificar este premio. El horizonte del evento
y el flujo contractual deben ser compatibles en una implementación empírica.

[Next verbatim source excerpt]

Proposición 1 (Solución global y pendientes por régimen). Bajo los Supuestos 1–3, existe un
único máximo en (2). Es
[U+001A control character]
[U+001B control character]
[U+001A control character]
[U+001A control character]
[U+001B control character][U+001B control character]
gR(D)
q 0 (D)
gR(D)
∗
0
∗
e = mı́n
, ē,
,
q = máx 0, q (D) − g mı́n ē,
.
(3)
c
g
c
Para κ fijo, en cualquier intervalo abierto con desigualdades estrictas que mantengan el régimen,
las pendientes son
[U+F8F1 delimiter glyph]
[U+F8F4 delimiter glyph]a − g 2 /c, gR/c < ē y gR/c < q 0 /g,
∂q ∗ [U+F8F2 delimiter glyph]
= a,
(4)
ē < gR/c y ē < q 0 /g,
[U+F8F4 delimiter glyph]
∂D
[U+F8F3 delimiter glyph]
0
0
0,
q /g < gR/c y q /g < ē.
Si ē = 0, e∗ = 0 y q ∗ = q 0 , con pendiente a en el interior del dominio de D. En empates de
restricciones la solución es continua y se usan derivadas laterales; una derivada ordinaria puede
no existir.

[Next verbatim source excerpt]

Fijemos un horizonte h. B0 > 0 es el saldo de créditos antiguos todavía vigente en esa fecha y
m0 ∈ [0, 1] su proporción en el estado de atraso elegido. B(D) > 0 es el saldo de la cohorte nueva
que permanece en cartera a h, después de amortización normal. Denote w(D) = B(D)/(B0 +B(D)).

[Next verbatim source excerpt]

Supuesto 5 (Contabilidad y maduración). Al horizonte, el saldo de cada cohorte es determinista
o se trata como exposición de una población continua. El riesgo q ∗ (D) coincide con la fracción
esperada de su saldo en el evento crediticio elegido, condicional a seguir en cartera. La amortización
y el tamaño del crédito no seleccionan riesgo adicional no recogido en q ∗ . B0 , m0 son constantes
al variar D; B(D) > 0 es diferenciable. No hay castigos, ventas de cartera ni reclasificaciones que
alteren esta identidad en el ejercicio.
\end{ReviewVerbatim}



\paragraph{Expanded PaperInterface specification}
\begin{ReviewVerbatim}
∀ (qbar a g c Rbar F E lower upper B0 m0 : ℝ) (B : ℝ → ℝ),
  Galarza26RelativeCredit.selectionAndDomain qbar a lower upper →
    Galarza26RelativeCredit.technologyCapacityCost g c E →
      Galarza26RelativeCredit.conditionedPayoff Rbar lower upper →
        Galarza26RelativeCredit.stockAccountingConditions B0 m0 lower upper B →
          (∀ D ∈ Set.Icc lower upper,
              Galarza26RelativeCredit.portfolioRatio B0 m0 B (Galarza26RelativeCredit.optimalRisk qbar a g c Rbar E) D =
                (1 - Galarza26RelativeCredit.cohortWeight B0 B D) * m0 +
                  Galarza26RelativeCredit.cohortWeight B0 B D *
                    Galarza26RelativeCredit.optimalRisk qbar a g c Rbar E D) ∧
            (∀ D ∈ Set.Icc lower upper,
                ∀ (beta Bprime : ℝ),
                  HasDerivAt (Galarza26RelativeCredit.optimalRisk qbar a g c Rbar E) beta D →
                    HasDerivAt B Bprime D →
                      have w : ℝ := Galarza26RelativeCredit.cohortWeight B0 B D;
                      have wprime : ℝ := Bprime * B0 / (B0 + B D) ^ 2;
                      have slope : ℝ :=
                        w * beta + wprime * (Galarza26RelativeCredit.optimalRisk qbar a g c Rbar E D - m0);
                      HasDerivAt
                          (Galarza26RelativeCredit.portfolioRatio B0 m0 B
                            (Galarza26RelativeCredit.optimalRisk qbar a g c Rbar E))
                          slope D ∧
                        ((∃ (b : ℝ), B = fun (x : ℝ) => b) →
                            (0 < slope ↔ 0 < beta) ∧ (slope = 0 ↔ beta = 0) ∧ (slope < 0 ↔ beta < 0)) ∧
                          (beta ≠ 0 →
                            |wprime * (Galarza26RelativeCredit.optimalRisk qbar a g c Rbar E D - m0)| < w * |beta| →
                              0 < slope * beta)) ∧
              ∀ D ∈ Set.Icc lower upper,
                ∀ (beta B1 epsilon : ℝ),
                  HasDerivAt (Galarza26RelativeCredit.optimalRisk qbar a g c Rbar E) beta D →
                    B = Galarza26RelativeCredit.exponentialBalance B1 epsilon →
                      0 < B1 →
                        0 ≤ epsilon →
                          have exponentialB : ℝ → ℝ := Galarza26RelativeCredit.exponentialBalance B1 epsilon;
                          have w : ℝ := Galarza26RelativeCredit.cohortWeight B0 exponentialB D;
                          have slope : ℝ :=
                            w *
                              (beta -
                                epsilon * (1 - w) * (Galarza26RelativeCredit.optimalRisk qbar a g c Rbar E D - m0));
                          HasDerivAt
                              (Galarza26RelativeCredit.portfolioRatio B0 m0 B
                                (Galarza26RelativeCredit.optimalRisk qbar a g c Rbar E))
                              slope D ∧
                            (0 < beta →
                              m0 < Galarza26RelativeCredit.optimalRisk qbar a g c Rbar E D →
                                (slope < 0 ↔
                                  beta <
                                    epsilon * (1 - w) * (Galarza26RelativeCredit.optimalRisk qbar a g c Rbar E D - m0)))
\end{ReviewVerbatim}

\paragraph{Retained governing declarations}
\begin{itemize}\raggedright
\item \hyperlink{paper-prerequisite-b0571b2c4a78eab5}{Galarza26\allowbreak{}Relative\allowbreak{}Credit.\allowbreak{}cohort\allowbreak{}Weight}
\item \hyperlink{paper-prerequisite-0bf7e0c122a51179}{Galarza26\allowbreak{}Relative\allowbreak{}Credit.\allowbreak{}conditioned\allowbreak{}Payoff}
\item \hyperlink{governing-declaration-ac78c6a99f925061}{Galarza26\allowbreak{}Relative\allowbreak{}Credit.\allowbreak{}exponential\allowbreak{}Balance}
\item \hyperlink{governing-declaration-e139cf967ddac229}{Galarza26\allowbreak{}Relative\allowbreak{}Credit.\allowbreak{}optimal\allowbreak{}Risk}
\item \hyperlink{governing-declaration-c80cb553f94a2c2b}{Galarza26\allowbreak{}Relative\allowbreak{}Credit.\allowbreak{}portfolio\allowbreak{}Ratio}
\item \hyperlink{paper-prerequisite-ab0e2f9b84884bc6}{Galarza26\allowbreak{}Relative\allowbreak{}Credit.\allowbreak{}selection\allowbreak{}And\allowbreak{}Domain}
\item \hyperlink{paper-prerequisite-de4cba6057d51d83}{Galarza26\allowbreak{}Relative\allowbreak{}Credit.\allowbreak{}stock\allowbreak{}Accounting\allowbreak{}Conditions}
\item \hyperlink{paper-prerequisite-e5f9b7ce782f743a}{Galarza26\allowbreak{}Relative\allowbreak{}Credit.\allowbreak{}technology\allowbreak{}Capacity\allowbreak{}Cost}
\end{itemize}
\paragraph{Lean-elaborated claim roles}\begin{ReviewVerbatim}
1. parameter: ℝ
2. parameter: ℝ
3. parameter: ℝ
4. parameter: ℝ
5. parameter: ℝ
6. parameter: ℝ
7. parameter: ℝ
8. parameter: ℝ
9. parameter: ℝ
10. parameter: ℝ
11. parameter: ℝ
12. parameter: ℝ → ℝ
13. assumption: Galarza26RelativeCredit.selectionAndDomain qbar a lower upper
14. assumption: Galarza26RelativeCredit.technologyCapacityCost g c E
15. assumption: Galarza26RelativeCredit.conditionedPayoff Rbar lower upper
16. assumption: Galarza26RelativeCredit.stockAccountingConditions B0 m0 lower upper B
17. conclusion: (∀ D ∈ Set.Icc lower upper,
    Galarza26RelativeCredit.portfolioRatio B0 m0 B (Galarza26RelativeCredit.optimalRisk qbar a g c Rbar E) D =
      (1 - Galarza26RelativeCredit.cohortWeight B0 B D) * m0 +
        Galarza26RelativeCredit.cohortWeight B0 B D * Galarza26RelativeCredit.optimalRisk qbar a g c Rbar E D) ∧
  (∀ D ∈ Set.Icc lower upper,
      ∀ (beta Bprime : ℝ),
        HasDerivAt (Galarza26RelativeCredit.optimalRisk qbar a g c Rbar E) beta D →
          HasDerivAt B Bprime D →
            have w : ℝ := Galarza26RelativeCredit.cohortWeight B0 B D;
            have wprime : ℝ := Bprime * B0 / (B0 + B D) ^ 2;
            have slope : ℝ := w * beta + wprime * (Galarza26RelativeCredit.optimalRisk qbar a g c Rbar E D - m0);
            HasDerivAt
                (Galarza26RelativeCredit.portfolioRatio B0 m0 B (Galarza26RelativeCredit.optimalRisk qbar a g c Rbar E))
                slope D ∧
              ((∃ (b : ℝ), B = fun (x : ℝ) => b) →
                  (0 < slope ↔ 0 < beta) ∧ (slope = 0 ↔ beta = 0) ∧ (slope < 0 ↔ beta < 0)) ∧
                (beta ≠ 0 →
                  |wprime * (Galarza26RelativeCredit.optimalRisk qbar a g c Rbar E D - m0)| < w * |beta| →
                    0 < slope * beta)) ∧
    ∀ D ∈ Set.Icc lower upper,
      ∀ (beta B1 epsilon : ℝ),
        HasDerivAt (Galarza26RelativeCredit.optimalRisk qbar a g c Rbar E) beta D →
          B = Galarza26RelativeCredit.exponentialBalance B1 epsilon →
            0 < B1 →
              0 ≤ epsilon →
                have exponentialB : ℝ → ℝ := Galarza26RelativeCredit.exponentialBalance B1 epsilon;
                have w : ℝ := Galarza26RelativeCredit.cohortWeight B0 exponentialB D;
                have slope : ℝ :=
                  w * (beta - epsilon * (1 - w) * (Galarza26RelativeCredit.optimalRisk qbar a g c Rbar E D - m0));
                HasDerivAt
                    (Galarza26RelativeCredit.portfolioRatio B0 m0 B
                      (Galarza26RelativeCredit.optimalRisk qbar a g c Rbar E))
                    slope D ∧
                  (0 < beta →
                    m0 < Galarza26RelativeCredit.optimalRisk qbar a g c Rbar E D →
                      (slope < 0 ↔
                        beta < epsilon * (1 - w) * (Galarza26RelativeCredit.optimalRisk qbar a g c Rbar E D - m0)))
\end{ReviewVerbatim}

\paragraph{Lean proof endpoint (built separately)}\begin{ReviewVerbatim}
Galarza26RelativeCredit.proposition3
\end{ReviewVerbatim}

\paragraph{Recorded source-to-Spec screening}
\textbf{Verdict:} \texttt{matches}
\\
{\raggedright
\textbf{Reason:} Compared the verbatim Proposition 3 and Assumptions 1–3 and 5 with the expanded target and all displayed accounting prerequisites.\allowbreak{} impairedBalance, portfolioRatio, and cohortWeight expand to (B0*m0+B*q*)/\allowbreak{}(B0+B) and w=B/\allowbreak{}(B0+B), so the first identity is exactly M=(1-\allowbreak{}w)m0+wq*.\allowbreak{} Under B0>0, m0 in [0,1], and positive differentiable B, the general derivative uses w'=B'B0/\allowbreak{}(B0+B)\textasciicircum{}2 and exactly M'=w*beta+w'*(q*-\allowbreak{}m0).\allowbreak{} The constant-\allowbreak{}B sign equivalences and the strict absolute-\allowbreak{}value dominance premise preserve the source's two sign rules.\allowbreak{} With B=B1*exp(-\allowbreak{}epsilon*D), B1>0, and epsilon>=0, the displayed slope is w\{beta-\allowbreak{}epsilon(1-\allowbreak{}w)(q*-\allowbreak{}m0)\}, and for beta>0 and q*>m0 its negative-\allowbreak{}slope equivalence is precisely epsilon(1-\allowbreak{}w)(q*-\allowbreak{}m0)>beta.\allowbreak{} No source case, hypothesis, formula, or conclusion is changed.\allowbreak{}
\par}
\reviewmatch{Matches source input}
\noindent\textbf{Reviewer annotation}\par
\reviewerbox

\end{document}
